@section Dialogs
The @code{Dialog} class implements all the functionality which is common
to the @code{ModalDialog} and @code{ModelessDialog} classes, and although
this class is not used directly, most of their functionality is documented
in this section.

@subsection  Standard Dialog Method Arguments

Since all dialog instance methods have the dialog object as the first
argument (referred to as @code{dlg}), this argument will not be described
each time.  It always refers to the dialog object which you wish to
perform the desired operation on.

Many of the methods associated with this class take an argument
identified as @code{id}.  This is a macro defined by the programmer via
the resource editor when the dialog is being defined and is used to
uniquely identify a particular control within the dialog.  Due to the
fact that this argument has the same meaning for every method which
uses it, it will not be defined each time.

The variable @code{ctl} will always refer to a control object.  Each
control object represents a unique control within a dialog.  Each
control object will be an instance of one of the subclasses of the
@code{Control} class (such as @code{TextControl} or
@code{NumericControl}).



@subsection Dialog Methods






@deffn {AddActivateFunctionAfter} AddActivateFunctionAfter::Dialog
@sp 2
@example
@group
r = gAddActivateFunctionAfter(dlg, fun);

object  dlg;        /*  a dialog object   */
int     (*fun)();   /*  function pointer  */
object  r;          /*  the dialog obj    */
@end group
@end example
This method is used to add a function which will be called when the dialog
is activated (receives focus).  The new function is added after any existing
activation handlers.

The value returned is the dialog object passed.
@example
@group
@exdent Example:

static int  my_activate(object dlg)
@{
        /*  handle activation  */
        return 0;
@}

gAddActivateFunctionAfter(dlg, my_activate);
@end group
@end example
@sp 1
See also:  @code{AddActivateFunctionBefore}
@end deffn




@deffn {AddActivateFunctionBefore} AddActivateFunctionBefore::Dialog
@sp 2
@example
@group
r = gAddActivateFunctionBefore(dlg, fun);

object  dlg;        /*  a dialog object   */
int     (*fun)();   /*  function pointer  */
object  r;          /*  the dialog obj    */
@end group
@end example
This method is used to add a function which will be called when the dialog
is activated (receives focus).  The new function is added before any existing
activation handlers.

The value returned is the dialog object passed.
@example
@group
@exdent Example:

static int  my_activate(object dlg)
@{
        /*  handle activation  */
        return 0;
@}

gAddActivateFunctionBefore(dlg, my_activate);
@end group
@end example
@sp 1
See also:  @code{AddActivateFunctionAfter}
@end deffn




@deffn {AddCompletionFunctionAfter} AddCompletionFunctionAfter::Dialog
@sp 2
@example
@group
r = gAddCompletionFunctionAfter(dlg, fun);

object  dlg;        /*  a dialog object   */
int     (*fun)();   /*  function pointer  */
object  r;          /*  the dialog obj    */
@end group
@end example
This method is used to add a function which will be called when the dialog
completes (when the user accepts or cancels the dialog).  The new function
is added after any existing completion handlers.

This is similar to @code{CompletionFunction} but allows multiple handlers
and controls the execution order.

The value returned is the dialog object passed.
@example
@group
@exdent Example:

static int  my_completion(object dlg)
@{
        /*  process completion  */
        return 0;
@}

gAddCompletionFunctionAfter(dlg, my_completion);
@end group
@end example
@sp 1
See also:  @code{AddCompletionFunctionBefore, CompletionFunction}
@end deffn




@deffn {AddCompletionFunctionBefore} AddCompletionFunctionBefore::Dialog
@sp 2
@example
@group
r = gAddCompletionFunctionBefore(dlg, fun);

object  dlg;        /*  a dialog object   */
int     (*fun)();   /*  function pointer  */
object  r;          /*  the dialog obj    */
@end group
@end example
This method is used to add a function which will be called when the dialog
completes (when the user accepts or cancels the dialog).  The new function
is added before any existing completion handlers.

This is similar to @code{CompletionFunction} but allows multiple handlers
and controls the execution order.

The value returned is the dialog object passed.
@example
@group
@exdent Example:

static int  my_completion(object dlg)
@{
        /*  process completion  */
        return 0;
@}

gAddCompletionFunctionBefore(dlg, my_completion);
@end group
@end example
@sp 1
See also:  @code{AddCompletionFunctionAfter, CompletionFunction}
@end deffn




@deffn {AddControl} AddControl::Dialog
@sp 2
@example
@group
ctl = mAddControl(dlg, ctlClass, id);

object  dlg;       /*  a dialog object   */
object  ctlClass;  /*  class of control  */
unsigned  id;      /*  control id        */
object  ctl;       /*  control object    */
@end group
@end example
This method is used to create a new control object and associate it to
a control defined in the dialog.  @code{ctlClass} must be literally
one of the subclasses of the @code{Control} class (such as
@code{TextControl}, @code{CheckBox}, @code{RadioButton}, etc).
See the class hierarchy for a complete list.

@code{id} identifies which control within dialog @code{dlg} to associate
the new control object with.

The value returned (@code{ctl}) is the new control object created and
will be an instance of the class indicated by the @code{ctlClass}
argument.  This object can be configured and queried via the protocol
defined by its class.  See that section for available options.  Setting
various options associated with this object will control default values
associated with the control and how this control functions within the
dialog.
@example
@group
@exdent Example:

object  dlg, ctl;

ctl = mAddControl(dlg, TextControl, PERSONS_NAME);
@end group
@end example
@sp 1
See also:  @code{GetControl}
@end deffn






@deffn {AddDlgHandlerAfter} AddDlgHandlerAfter::Dialog
@sp 2
@example
@group
r = gAddDlgHandlerAfter(dlg, msg, func);

object   dlg;      /*  a dialog object  */
unsigned msg;      /*  message          */
BOOL    (*func)(); /*  function pointer */
object  r;         /*  the dialog obj   */
@end group
@end example
This method is used to associate function @code{func} with Windows dialog
message @code{msg} for dialog @code{dlg}.  Whenever dialog @code{dlg}
receives message @code{msg}, @code{func} will be called.

@code{dlg} is the dialog object who's messages you wish to process.
@code{msg} is the particular message you wish to trap.  These messages
are fully documented in the Windows documentation in the Messages
section.  They normally begin with @code{WM_}.

@code{func} is the function which gets called whenever the specified
message gets received and takes the following form:
@example
@group
BOOL    func(object     dlg,
             HWND       hdlg, 
             UINT       mMsg, 
             WPARAM     wParam, 
             LPARAM     lParam)
@{
        .
        .
        .
        return FALSE;  /* or whatever is appropriate  */
@}
@end group
@end example
Where @code{dlg} is the dialog being sent the message.  The remaining
arguments and return value is fully documented in the Windows documentation
under the @code{DialogProc} function and the Windows Messages documentation.

WDS keeps a list of functions associated with each message associated
with each dialog.  When a particular message is received the appropriate
list of handler functions gets executed sequentially.
@code{AddDlgHandlerAfter} appends the new function to the end of this list,
and @code{AddDlgHandlerBefore} adds the new function to the beginning of
the list.

Windows will only see the return value of the last message handler executed.
@example
@group
@exdent Example:

int     hSize, vSize;

static  BOOL    process_wm_size(object  dlg, 
                                HWND    hwnd, 
                                UINT    mMsg, 
                                WPARAM  wParam, 
                                LPARAM  lParam)
@{
        hSize = LOWORD(lParam);
        vSize = HIWORD(lParam);
        return 0L;
@}

        .
        .
        gAddDlgHandlerAfter(dlg, (unsigned) WM_SIZE,
                                      process_wm_size);
        .
        .
@end group
@end example
@sp 1
See also:  @code{AddDlgHandlerBefore}
@end deffn












@deffn {AddDlgHandlerBefore} AddDlgHandlerBefore::Dialog
@sp 2
@example
@group
r = gAddDlgHandlerBefore(dlg, msg, func);

object   dlg;     /*  a dialog object  */
unsigned msg;      /*  message          */
BOOL    (*func)(); /*  function pointer */
object  r;         /*  the dialog obj   */
@end group
@end example
This function is fully documented under @code{AddDlgHandlerAfter}.
@c @example
@c @group
@c @exdent Example:
@c @end group
@c @end example
@sp 1
See also:  @code{AddDlgHandlerAfter}
@end deffn













@deffn {AddInitFunctionAfter} AddInitFunctionAfter::Dialog
@sp 2
@example
@group
r = gAddInitFunctionAfter(dlg, fun);

object  dlg;        /*  a dialog object   */
int     (*fun)();   /*  function pointer  */
object  r;          /*  the dialog obj    */
@end group
@end example
This method is used to add a function which will be called after the dialog
has been initialized.  The new function is added after any existing
initialization handlers.

The value returned is the dialog object passed.
@example
@group
@exdent Example:

static int  my_init(object dlg)
@{
        /*  perform initialization  */
        return 0;
@}

gAddInitFunctionAfter(dlg, my_init);
@end group
@end example
@sp 1
See also:  @code{AddInitFunctionBefore}
@end deffn




@deffn {AddInitFunctionBefore} AddInitFunctionBefore::Dialog
@sp 2
@example
@group
r = gAddInitFunctionBefore(dlg, fun);

object  dlg;        /*  a dialog object   */
int     (*fun)();   /*  function pointer  */
object  r;          /*  the dialog obj    */
@end group
@end example
This method is used to add a function which will be called before the dialog
has been initialized.  The new function is added before any existing
initialization handlers.

The value returned is the dialog object passed.
@example
@group
@exdent Example:

static int  my_init(object dlg)
@{
        /*  perform initialization  */
        return 0;
@}

gAddInitFunctionBefore(dlg, my_init);
@end group
@end example
@sp 1
See also:  @code{AddInitFunctionAfter}
@end deffn




@deffn {AutoDispose} AutoDispose::Dialog
@sp 2
@example
@group
pflg = gAutoDispose(dlg, flg);

object  dlg;    /*  a dialog object    */
int     flg;    /*  desired mode       */
int     pflg;   /*  previous mode set  */
@end group
@end example
This method is used to control the auto disposal facility associated with
a given dialog.  The default is disabled for modal dialogs and enabled
for modeless dialogs.

If enabled, the auto dispose feature will cause the dialog object
(@code{dlg}) to be automatically disposed of whenever the user accepts
or aborts the dialog.  If disabled, the application may continue to use
the object (for querying control values for example) after the dialog
has been terminated.  If the auto dispose feature is disabled the
application must explicitly dispose of the object when it is no longer
needed.

Note that if the auto dispose feature is used, the @code{gCompletionFunction}
method may be used to access the control values immediately prior to
the object's automatic disposal.

Set @code{flg} to 1 to enable the feature and 0 to disable it.  The
previous state will be returned.
@example
@group
@exdent Example:

object  dlg;

gAutoDispose(dlg, 1);
@end group
@end example
@sp 1
See also:  @code{CompletionFunction, Perform, SetResult}
@end deffn














@deffn {BackBrush} BackBrush::Dialog
@sp 2
@example
@group
r = gBackBrush(dlg, brsh);

object  dlg;    /*  a dialog object   */
object  brsh;   /*  brush object      */
object  r;      /*  dlg               */
@end group
@end example
This method is used to set the brush object which is used for the
background of the dialog.  Any previously associated brush object will
be disposed.  This brush object will also be automatically disposed when
the dialog is disposed.

The dialog passed is returned.
@example
@group
@exdent Example:

object  dlg;

gBackBrush(dlg, vNew(SolidBrush, 0, 255, 0));
@end group
@end example
@sp 1
See also:  @code{TextBrush, SetBackBrush::Application}
        and the @code{Brush} classes
@end deffn










@deffn {Center} Center::Dialog
@sp 2
@example
@group
r = gCenter(dlg);

object  dlg;   /*  a dialog object  */
object  r;     /*  the dialog obj   */
@end group
@end example
This method is used to center the dialog on its parent window.  If the
dialog is currently active, it is centered immediately.  If the dialog
is not yet visible, it will be centered when it is shown.

The value returned is the dialog object passed.
@example
@group
@exdent Example:

object  dlg;

gCenter(dlg);
@end group
@end example
@end deffn




@deffn {CheckValue} CheckValue::Dialog
@sp 2
@example
@group
r = gCheckValue(dlg);

object  dlg;   /*  a dialog object   */
int     r;     /*  validation result  */
@end group
@end example
This method is used to validate all controls in the dialog.  Each control's
check function (if any) is called to verify its current value.

The value returned is 0 if all controls pass validation, or 1 if any
control has a validation error.
@example
@group
@exdent Example:

object  dlg;

if (gCheckValue(dlg))
        /*  handle validation error  */
@end group
@end example
@end deffn




@deffn {CompletionFunction} CompletionFunction::Dialog
@sp 2
@example
@group
r = gCompletionFunction(dlg, fun);

object  dlg;       /*  dialog object        */
int     (*fun)();  /*  completion function  */
object  r;         /*  dialog object        */
@end group
@end example
This method is used to associate a C function with a dialog such that
when the user accepts or aborts the dialog, function @code{fun} will
get executed by WDS.  The function will get executed prior to any automatic
disposal of the dialog object.

The completion function facility is best used in conjunction with
modeless dialogs and the auto dispose facility in order to process
control values when the user accepts the dialog and prior to the
automatic disposal of the dialog object.

The format of @code{fun} is as follows:
@example
@group
r = fun(dlg, res);

object  dlg;    /*  the dialog object  */
int     res;    /*  result status of the dialog  */
int     r;      /*  result for gPerform  */
@end group
@end example
Where @code{res} will be @code{TRUE} if the user accepted the dialog or
@code{FALSE} if the user canceled the dialog.  @code{r} is the value
which will become the result of the dialog which is returned by @code{gPerform}
(if it was a modal dialog).  @code{r} will normally be @code{res}.
@example
@group
@exdent Example:

static  int     fun(object dlg, int res)
@{
        .
        .
        return res;
@}

        gCompletionFunction(dlg, fun);
@end group
@end example
@sp 1
See also:  @code{AutoDispose, Perform, SetTag, SetResult}
@end deffn














@deffn {CtlDoubleValue} CtlDoubleValue::Dialog
@sp 2
@example
@group
val = mCtlDoubleValue(dlg, id);

object  dlg;    /*  a dialog object  */
unsigned id;    /*  control id       */
double  val;    /*  control value    */
@end group
@end example
This method is used to gain access to the value associated with a
particular control as a C double.  For example, if the control
was a numeric control and the user entered ``3.141'', then this method
would return a double ``3.141''.
@example
@group
@exdent Example:

object  dlg;
double  val;

val = mCtlDoubleValue(dlg, PI_VALUE);
@end group
@end example
@sp 1
See also:  @code{CtlStringValue, CtlShortValue, ...}
        @code{CtlValue}
@end deffn

















@deffn {CtlLongValue} CtlLongValue::Dialog
@sp 2
@example
@group
val = mCtlLongValue(dlg, id);

object  dlg;    /*  a dialog object  */
unsigned id;    /*  control id       */
long    val;    /*  control value    */
@end group
@end example
This method is used to gain access to the value associated with a
particular control as a C long.  For example, if the control
was a numeric control and the user entered ``6'', then this method
would return a long ``6''.

In the case of date controls, the value returned will be in the
form YYYYMMDD.  For example November 24, 1956 would be
19561124.
@example
@group
@exdent Example:

object  dlg;
long    val;

val = mCtlLongValue(dlg, PERSONS_AGE);
@end group
@end example
@sp 1
See also:  @code{CtlStringValue, CtlShortValue, ...}
        @code{CtlValue}
@end deffn














@deffn {CtlShortValue} CtlShortValue::Dialog
@sp 2
@example
@group
val = mCtlShortValue(dlg, id);

object  dlg;    /*  a dialog object  */
unsigned id;    /*  control id       */
short   val;    /*  control value    */
@end group
@end example
This method is used to gain access to the value associated with a
particular control as a C short.  For example, if the control
was a numeric control and the user entered ``6'', then this method
would return a short ``6''.

In the case of a list box or combo box the value returned will be in
terms of an ordinal value. This value returned is a zero based index
from top to bottom indicating the selection made by the user.  If no
selection was made negative value will be returned.  

In the case of check boxes and radio buttons this method will return
0 if the button is not selected, 1 if the button is selected, and
2 if the button is grayed.

In the case of scroll bars, this method will return a short value
between the minimum and maximum set for the control which indicates
where the control was set to.

@code{mCtlLongValue} should be used for date controls.
@example
@group
@exdent Example:

object  dlg;
short   val;

val = mCtlShortValue(dlg, PERSONS_AGE);
@end group
@end example
@sp 1
See also:  @code{CtlStringValue, CtlLongValue, ...}
        @code{CtlValue}
@end deffn















@deffn {CtlStringValue} CtlStringValue::Dialog
@sp 2
@example
@group
val = mCtlStringValue(dlg, id);

object  dlg;    /*  a dialog object  */
unsigned id;    /*  control id       */
char    *val;   /*  control value    */
@end group
@end example
This method is used to gain access to the value associated with a
particular control as a C character string.  For example, if the control
was a text control and the user entered ``Miami'', then this method
would return a pointer to the C string ``Miami''.

If the control being accessed is a list box or combo box, the value
returned will be a string representation of the text associated with the
choice made by the user.  See @code{CtlShortValue} or @code{IndexValue}
to get an ordinal value for this control type.

The character pointer returned will not be valid once the dialog or
control has been disposed.  Therefore, the value should be kept by
some other means if it is desired past the life of the dialog.
@example
@group
@exdent Example:

object  dlg;
char    *val;

val = mCtlStringValue(dlg, PERSONS_NAME);
@end group
@end example
@sp 1
See also:  @code{CtlShortValue, CtlLongValue, ...}
        @code{CtlValue}
@end deffn











@deffn {CtlUnsignedShortValue} CtlUnsignedShortValue::Dialog
@sp 2
@example
@group
val = mCtlUnsignedShortValue(dlg, id);

object  dlg;    /*  a dialog object  */
unsigned id;    /*  control id       */
unsigned short  val;    /*  control value    */
@end group
@end example
This method is used to gain access to the value associated with a
particular control as a C unsigned short.  For example, if the control
was a numeric control and the user entered ``6'', then this method
would return a unsigned short ``6''.

@code{mCtlLongValue} should be used for date controls.
@example
@group
@exdent Example:

object  dlg;
unsigned short  val;

val = mCtlUnsignedShortValue(dlg, PERSONS_AGE);
@end group
@end example
@sp 1
See also:  @code{CtlStringValue, CtlLongValue, ...}
        @code{CtlValue}
@end deffn














@deffn {CtlValue} CtlValue::Dialog
@sp 2
@example
@group
val = mCtlValue(dlg, id);

object  dlg;    /*  a dialog object  */
unsigned id;    /*  control id       */
object  val;    /*  control value    */
@end group
@end example
This method is used to gain access to the value associated with a
particular control as a Dynace object.  For example, if the control
was a text control and user entered ``Miami'', then this method
would return a Dynace string object (not a normal C string) which
represents the value typed in.  This returned value may be used
and accessed like any other Dynace object using the appropriate
interface mechanism.

If the control being accessed is a list box or combo box, the value
returned will be a representation of the text associated with the choice
made by the user.  See @code{IndexValue} or @code{CtlShortValue} to get
a numeric representation for this control type.

The object returned will be automatically disposed when the dialog is
disposed, therefore, a copy (via @code{gCopy}) should be made of it if
its existence is required subsequent to the dialog's disposal.

WDS also provides several methods to obtain normal C data types directly.
@example
@group
@exdent Example:

object  dlg, val;

val = mCtlValue(dlg, PERSONS_NAME);
@end group
@end example
@sp 1
See also:  @code{CtlStringValue, CtlShortValue, CtlLongValue, ...}
@end deffn












@deffn {DeepDispose} DeepDispose::Dialog
@sp 2
This method performs the same function as @code{Dispose}.  See that
method for details.
@end deffn









@deffn {Disable} Disable::Dialog
@sp 2
@example
@group
r = gDisable(dlg);

object  dlg;   /*  a dialog object  */
object  r;     /*  the dialog obj   */
@end group
@end example
This method is used to disable the dialog window, preventing user input.
A disabled dialog will not respond to mouse or keyboard events until
it is re-enabled with @code{gEnable}.

The value returned is the dialog object passed.
@example
@group
@exdent Example:

object  dlg;

gDisable(dlg);
@end group
@end example
@sp 1
See also:  @code{Enable}
@end deffn




@deffn {DisableAll} DisableAll::Dialog
@sp 2
@example
@group
r = gDisableAll(dlg);

object  dlg;   /*  a dialog object  */
object  r;     /*  the dialog obj   */
@end group
@end example
This method is used to disable all controls contained within the dialog,
preventing user input to any of them.  Individual controls or all controls
can be re-enabled with @code{gEnableAll}.

The value returned is the dialog object passed.
@example
@group
@exdent Example:

object  dlg;

gDisableAll(dlg);
@end group
@end example
@sp 1
See also:  @code{EnableAll}
@end deffn




@deffn {Dispose} Dispose::Dialog
@sp 2
@example
@group
r = gDispose(dlg);

object  dlg;   /*  a dialog object   */
object  r;     /*  NULL              */
@end group
@end example
This method is used to remove and dispose of a dialog object when it
is no longer needed.  This method should be called on all dialogs
when they are no longer needed.  In addition to disposing of the dialog,
this method will also dispose of all the control objects associated
with the dialog.

The value returned is always @code{NULL} and may be used to null out
the variable which contained the object being disposed in order to
avoid future accidental use.
@example
@group
@exdent Example:

object  dlg;

dlg = gDispose(dlg);
@end group
@end example
@sp 1
See also:  @code{AutoDispose}
@end deffn










@deffn {Enable} Enable::Dialog
@sp 2
@example
@group
r = gEnable(dlg);

object  dlg;   /*  a dialog object  */
object  r;     /*  the dialog obj   */
@end group
@end example
This method is used to enable the dialog window for user input.  A dialog
that has been previously disabled with @code{gDisable} can be re-enabled
with this method.

The value returned is the dialog object passed.
@example
@group
@exdent Example:

object  dlg;

gEnable(dlg);
@end group
@end example
@sp 1
See also:  @code{Disable}
@end deffn




@deffn {EnableAll} EnableAll::Dialog
@sp 2
@example
@group
r = gEnableAll(dlg);

object  dlg;   /*  a dialog object  */
object  r;     /*  the dialog obj   */
@end group
@end example
This method is used to enable all controls contained within the dialog.
This is useful for re-enabling all controls after a previous call to
@code{gDisableAll}.

The value returned is the dialog object passed.
@example
@group
@exdent Example:

object  dlg;

gEnableAll(dlg);
@end group
@end example
@sp 1
See also:  @code{DisableAll}
@end deffn




@deffn {GetBackBrush} GetBackBrush::Dialog
@sp 2
@example
@group
brsh = gGetBackBrush(dlg);

object  dlg;    /*  dialog object  */
object  brsh;   /*  brush object   */
@end group
@end example
This method is used to obtain the background brush object which has been
associated with the dialog.  If no specific background brush was
associated with the dialog this method will return the application
default background brush, which is what the dialog uses when no specific
brush is specified.
@example
@group
@exdent Example:

object  dlg, brsh;

brsh = gGetBackBrush(dlg);
@end group
@end example
@sp 1
See also:  @code{BackBrush, GetTextBrush, GetBackBrush::Application}
@end deffn









@deffn {GetControl} GetControl::Dialog
@sp 2
@example
@group
ctl = mGetControl(dlg, id);

object  dlg;    /*  a dialog object  */
unsigned id;    /*  control id       */
object  ctl;    /*  control object   */
@end group
@end example
This method is used to gain access to a control object which has been
previously associated with a particular control via its associated
control id.  The control object is returned.
@example
@group
@exdent Example:

object  dlg, ctl;

ctl = mGetControl(dlg, PERSONS_NAME);
@end group
@end example
@sp 1
See also:  @code{AddControl}
@end deffn












@deffn {GetParent} GetParent::Dialog
@sp 2
@example
@group
prnt = gGetParent(dlg);

object  dlg;    /*  dialog object         */
object  prnt;   /*  parent window object  */
@end group
@end example
This method is used to obtain the parent window object associated
with dialog @code{dlg}.  This parent window would have established
when the dialog object was created.
@example
@group
@exdent Example:

object  dlg, parentWind;

parentWind = gGetParent(dlg);
@end group
@end example
@sp 1
See also:  @code{NewDialog::ModalDialog, NewDialog::ModelessDialog}
@end deffn















@deffn {GetPosition} GetPosition::Dialog
@sp 2
@example
@group
r = gGetPosition(dlg, y, x);

object  dlg;   /*  a dialog object     */
int     *y;    /*  vertical position   */
int     *x;    /*  horizontal position */
object  r;     /*  the dialog obj      */
@end group
@end example
This method is used to obtain the current screen position of the dialog.
The position is returned in the current scaling mode through the
@code{y} and @code{x} pointer arguments.

The value returned is the dialog object passed.
@example
@group
@exdent Example:

object  dlg;
int     yPos, xPos;

gGetPosition(dlg, &yPos, &xPos);
@end group
@end example
@sp 1
See also:  @code{GetSize}
@end deffn




@deffn {GetResourceID} GetResourceID::Dialog
@sp 2
@example
@group
r = gGetResourceID(dlg);

object    dlg;   /*  a dialog object   */
unsigned  r;     /*  resource ID       */
@end group
@end example
This method is used to obtain the resource identifier associated with the
dialog.  The resource ID is the value used when the dialog was created.

The value returned is the dialog's resource identifier.
@example
@group
@exdent Example:

object    dlg;
unsigned  id;

id = gGetResourceID(dlg);
@end group
@end example
@end deffn




@deffn {GetSize} GetSize::Dialog
@sp 2
@example
@group
r = gGetSize(dlg, y, x);

object  dlg;   /*  a dialog object   */
int     *y;    /*  vertical size     */
int     *x;    /*  horizontal size   */
object  r;     /*  the dialog obj    */
@end group
@end example
This method is used to obtain the current size of the dialog.
The size is returned in the current scaling mode through the
@code{y} and @code{x} pointer arguments.

The value returned is the dialog object passed.
@example
@group
@exdent Example:

object  dlg;
int     height, width;

gGetSize(dlg, &height, &width);
@end group
@end example
@sp 1
See also:  @code{GetPosition}
@end deffn




@deffn {GetTag} GetTag::Dialog
@sp 2
@example
@group
r = gGetTag(dlg);

object  dlg;   /*  dialog object  */
object  r;     /*  tag            */
@end group
@end example
This method is used to obtain a Dynace object which has been associated
with a dialog via @code{SetTag}.  The value return is the object which has
been associated with the dialog object @code{dlg}.  If there is no object
associated with the dialog, @code{NULL} will be returned.
@example
@group
@exdent Example:

object  dlg, someObj;

someObj = gGetTag(dlg);
@end group
@end example
@sp 1
See also:  @code{SetTag, SetTag::Window}
@end deffn










@deffn {GetTask} GetTask::Dialog
@sp 2
@example
@group
r = gGetTask(dlg);

object  dlg;   /*  a dialog object  */
object  r;     /*  Task object      */
@end group
@end example
This method is used to obtain the Task object associated with the dialog.
If no Task has been associated, @code{NULL} is returned.

The value returned is the associated Task object or @code{NULL}.
@example
@group
@exdent Example:

object  dlg, task;

task = gGetTask(dlg);
@end group
@end example
@sp 1
See also:  @code{SetTask}
@end deffn




@deffn {GetTextBrush} GetTextBrush::Dialog
@sp 2
@example
@group
brsh = gGetTextBrush(dlg);

object  dlg;    /*  dialog object  */
object  brsh;   /*  brush object   */
@end group
@end example
This method is used to obtain the text brush object which has been associated
with the dialog.  If no specific text brush was associated with the dialog
this method will return the application default text brush, which is what
the dialog uses when no specific brush is specified.
@example
@group
@exdent Example:

object  dlg, brsh;

brsh = gGetTextBrush(dlg);
@end group
@end example
@sp 1
See also:  @code{TextBrush, GetBackBrush, GetTextBrush::Application}
@end deffn











@deffn {GetTopic} GetTopic::Dialog
@sp 2
@example
@group
r = gGetTopic(dlg);

object  dlg;   /*  a dialog object   */
char    *r;    /*  help topic string */
@end group
@end example
This method is used to obtain the help topic string associated with the
dialog.  If no topic has been set, @code{NULL} is returned.

The value returned is a pointer to the help topic string or @code{NULL}.
@example
@group
@exdent Example:

object  dlg;
char    *topic;

topic = gGetTopic(dlg);
@end group
@end example
@sp 1
See also:  @code{SetTopic}
@end deffn




@deffn {Handle} Handle::Dialog
@sp 2
@example
@group
h = gHandle(dlg);

object  dlg;    /*  dialog object   */
HANDLE  h;      /*  Windows handle  */
@end group
@end example
This method is used to obtain the Windows internal handle associated with
a dialog object.  It will only return a valid handle while the dialog
is being performed via @code{gPerform}.  @code{NULL} will be returned
otherwise.

Note that this method may be used with most WDS objects in order to obtain
the internal handle that Windows normally associates with each type of object.
See the appropriate documentation.
@example
@group
@exdent Example:

object  dlg;
HANDLE  h;

h = gHandle(dlg);
@end group
@end example
@c @sp 1
@c See also:  @code{Message}
@end deffn








@deffn {IndexValue} IndexValue::Dialog
@sp 2
@example
@group
val = mIndexValue(dlg, id);

object  dlg;    /*  a dialog object  */
unsigned id;    /*  control id       */
object  val;    /*  control value    */
@end group
@end example
This method is used to gain access to a Dynace object (as opposed to a
normal C integer) representing the value associated with a list box or
combo box in terms of an ordinal value. This value returned is a zero
based index from top to bottom indicating the selection made by the
user.  If no selection was made the object will represent a negative
value.  This returned value may be used and accessed like any other
Dynace object using the appropriate interface mechanisms.

The object returned will be automatically disposed when the dialog is
disposed, therefore, a copy (via @code{gCopy}) should be made of it if
its existence is required subsequent to the dialog's disposal.

WDS also provides several methods to obtain normal C data types directly.
@example
@group
@exdent Example:

object  dlg, val;

val = mIndexValue(dlg, SOME_LISTBOX);
@end group
@end example
@sp 1
See also:  @code{CtlShortValue, CtlValue}
@end deffn








@deffn {InDialog} InDialog::Dialog
@sp 2
@example
@group
flg = gInDialog(dlg);

object  dlg;    /*  a dialog object  */
int     flg;    /*  in dialog flag   */
@end group
@end example
This method is used to determine whether a given dialog is currently
being ``performed'' via @code{gPerform}.  Prior and subsequent to
performing a dialog this method will return 0.  A 1 will be returned
if the dialog is currently being performed.
@example
@group
@exdent Example:

object  dlg;
int     flg;

flg = gInDialog(dlg);
@end group
@end example
@sp 1
See also:  @code{Perform}
@end deffn















@deffn {Message} Message::Dialog
@sp 2
@example
@group
r = gMessage(dlg, msg);

object  dlg;    /*  dialog object  */
char    *msg;   /*  message        */
object  r;      /*  dialog object  */
@end group
@end example
This method is used to open up a temporary informational window.  The
window will contain the message given by @code{msg} and the user must
acknowledge the window prior to continuing by hitting an OK button.

This method may only be used either while the dialog is being performed
or at any time if it has an associated parent window.

The value returned is the dialog passed.
@example
@group
@exdent Example:

object  dlg;

gMessage(dlg, "Press OK to continue.");
@end group
@end example
@sp 1
See also:  @code{MessageWithTopic}
@end deffn












@deffn {MessageWithTopic} MessageWithTopic::Dialog
@sp 2
@example
@group
r = gMessageWithTopic(dlg, msg, tpc);

object  dlg;    /*  dialog object  */
char    *msg;   /*  message        */
char    *tpc;   /*  help topic     */
object  r;      /*  dialog object  */
@end group
@end example
This method is used to open up a temporary informational dialog.  The
dialog will contain the message given by @code{msg} and the user must
acknowledge the dialog prior to continuing by hitting an OK button.

If the user hits the @code{F1} key while presented with the message,
the help topic identified by @code{tpc} will get displayed via the Windows
help system.

This method may only be used either while the dialog is being performed
or at any time if it has an associated parent window.

The value returned is the dialog passed.
@example
@group
@exdent Example:

object  dlg;

gMessageWithTopic(dlg, "Press OK to continue.", "mytopic");
@end group
@end example
@sp 1
See also:  @code{Message} and the @code{HelpSystem} class.
@end deffn














@deffn {MoveToTop} MoveToTop::Dialog
@sp 2
@example
@group
r = gMoveToTop(dlg);

object  dlg;   /*  a dialog object  */
object  r;     /*  the dialog obj   */
@end group
@end example
This method is used to bring the dialog to the top of the Z-order,
making it the foremost window.  This is a convenience method equivalent
to calling @code{gSetZOrder} with @code{HWND_TOP}.

The value returned is the dialog object passed.
@example
@group
@exdent Example:

object  dlg;

gMoveToTop(dlg);
@end group
@end example
@sp 1
See also:  @code{SetZOrder}
@end deffn




@deffn {Perform} Perform::Dialog
@sp 2
@example
@group
r = gPerform(dlg);

object  dlg;    /*  a dialog object  */
int     r;      /*  result           */
@end group
@end example
Once a dialog is created and fully specified, this method is used to
actually display the dialog and allow the user to interact with it.
The dialog will remain active until the user pushes the push buttons
with the label @code{IDOK} or @code{IDCANCEL}.

What occurs after calling @code{gPerform} depends on whether the dialog
was a modal or modeless dialog.

Control initial values should be set prior to calling @code{gPerform}.

MODAL DIALOGS

If the dialog being performed is a modal, then once @code{gPerform} is
evoked the user will be able to interact with the dialog and @code{gPerform}
will not return until the user either accepts or aborts the entire dialog.
This, in effect, will prevent the user from switching to any other window
or dialog within the application until the dialog is completed.

Finally, the value returned by @code{gPerform} will be the Windows
constant @code{FALSE} if the user canceled the dialog, the Windows
constant @code{TRUE} if the user accepted the dialog, or the
value returned by any completion function which was attached to the
dialog (via @code{CompletionFunction}).

Once @code{gPerform} returns, the dialog object may be queried in
order to obtain the final values associated with each control.

Unless the auto dispose feature is manually set for a modal dialog,
the dialog object is not automatically disposed and must be when
no longer needed.

MODELESS DIALOGS

Modeless dialogs operate different from modal dialogs.  When
@code{gPerform} is evoked on a modeless dialog, @code{gPerform} returns
immediately with a 0 value.  At that point the user is presented with the
dialog and allowed to edit it.  During that time your program would typically
return to the menu and wait for further user selections.  The user may
continue with the current dialog, abort it, or switch to another window
within the application with the ability to return at any point.  All the
user interaction is automatically handled by Windows and WDS.

Since the application code has gone on about its business after a call
to @code{gPerform} two things are left to be completed whenever the user
finally accepts or aborts the dialog.  First, the application will
typically want to do something with the data the user has entered on the
dialog.  And second, the dialog object will need to be disposed when
it's no longer needed.

When using modeless dialogs the programmer normally attaches a completion
function to the dialog via @code{gCompletionFunction}.  What this does
is cause the completion function to be executed whenever the user
accepts or aborts the dialog.  This application specific function can be
used to perform the final processing of the dialog's data.  See
@code{gCompletionFunction} for complete details.

By default, the auto dispose flag associated with modeless dialogs is
normally enabled.  This causes WDS to automatically dispose of a dialog
object and all associated controls immediately subsequent to executing
the completion function.  This way everything is cleaned up automatically.
This feature may be disabled via @code{gAutoDispose}, in which case the
application must be sure to dispose of the unneeded object.
@example
@group
@exdent Example:

object  dlg;
int     r;

r = gPerform(dlg);
@end group
@end example
@sp 1
See also:  @code{CompletionFunction, AutoDispose}
@end deffn













@deffn {PressButton} PressButton::Dialog
@sp 2
@example
@group
r = gPressButton(dlg, id);

object    dlg;   /*  a dialog object   */
unsigned  id;    /*  button control id */
object    r;     /*  the dialog obj    */
@end group
@end example
This method is used to programmatically simulate clicking a button control
within the dialog.  The @code{id} argument identifies the button by its
resource id.  If no button with the given id is found, @code{NULL} is
returned.

The value returned is the dialog object passed, or @code{NULL} if the
button was not found.
@example
@group
@exdent Example:

object  dlg;

gPressButton(dlg, IDOK);
@end group
@end example
@end deffn




















@deffn {SetCursor} SetCursor::Dialog
@sp 2
@example
@group
r = gSetCursor(dlg, cursor);

object  dlg;      /*  a dialog object  */
object  cursor;   /*  cursor object    */
object  r;        /*  the dialog obj   */
@end group
@end example
This method is used to set the cursor for the dialog and all of its
controls.  Pass @code{NULL} for @code{cursor} to restore the default
cursor.

The value returned is the dialog object passed.
@example
@group
@exdent Example:

object  dlg, cur;

gSetCursor(dlg, cur);
gSetCursor(dlg, NULL);  /*  restore default  */
@end group
@end example
@sp 1
See also:  @code{WaitCursor}
@end deffn




@deffn {SetFocus} SetFocus::Dialog
@sp 2
@example
@group
r = gSetFocus(dlg);

object  dlg;   /*  a dialog object  */
object  r;     /*  the dialog obj   */
@end group
@end example
This method is used to set the keyboard input focus to the dialog window.
This is useful when the dialog needs to regain focus after another window
has been activated.

The value returned is the dialog object passed.
@example
@group
@exdent Example:

object  dlg;

gSetFocus(dlg);
@end group
@end example
@end deffn




@deffn {SetResult} SetResult::Dialog
@sp 2
@example
@group
r = gSetResult(dlg, obj);

object  dlg;    /*  a dialog object  */
object  obj;    /*  result object    */
object  r;      /*  dlg              */
@end group
@end example
This method is used to associate a Dynace short integer object with a
dialog such that when the dialog is completed or canceled by the user,
the associated object will be set to the dialog's result value.  This
feature is most often needed in conjunction with modeless dialogs and
the auto dispose feature.  It can be used to find out what happened to a
dialog subsequent to the dialog object being disposed.

The object associated with the dialog will never be disposed by WDS
and must be manually disposed when it is no longer needed.
@example
@group
@exdent Example:

object  dlg, obj;

gSetResult(dlg, obj = gNewWithInt(ShortInteger, -1));
@end group
@end example
@sp 1
See also:  @code{AutoDispose, Perform, SetTag}
@end deffn










@deffn {SetState} SetState::Dialog
@sp 2
@example
@group
r = gSetState(dlg, state);

object  dlg;     /*  a dialog object  */
int     state;   /*  show state       */
object  r;       /*  the dialog obj   */
@end group
@end example
This method is used to set the display state of the dialog window.
The @code{state} argument uses the standard Windows show state constants
such as @code{SW_SHOW}, @code{SW_HIDE}, @code{SW_MINIMIZE},
@code{SW_MAXIMIZE}, and @code{SW_RESTORE}.

The value returned is the dialog object passed.
@example
@group
@exdent Example:

object  dlg;

gSetState(dlg, SW_HIDE);
@end group
@end example
@end deffn




@deffn {SetTag} SetTag::Dialog
@sp 2
@example
@group
r = gSetTag(dlg, tag);

object  dlg;   /*  a dialog object   */
object  tag;    /*  tag               */
object  r;      /*  previous tag      */
@end group
@end example
This method is used to associate an arbitrary Dynace object with a
dialog object.  This may later be retrieved via the @code{GetTag} method.
Since WDS passes around the dialog object to all @code{Dialog} methods
this mechanism may be used to pass additional information with the
dialog.  And since Dynace treats all objects in a uniform manner, this
information attached to the dialog may be arbitrarily complex.

WDS does not dispose of the tag when the dialog object is disposed.
This method returns any previous object associated with the dialog or
@code{NULL}.
@example
@group
@exdent Example:

object  dlg;

gSetTag(dlg, gNewWithInt(ShortInteger, 17));
@end group
@end example
@sp 1
See also:  @code{GetTag, SetTag::Window}
@end deffn











@deffn {SetTask} SetTask::Dialog
@sp 2
@example
@group
r = gSetTask(dlg, task);

object  dlg;    /*  a dialog object  */
object  task;   /*  Task object      */
object  r;      /*  the dialog obj   */
@end group
@end example
This method is used to associate a Task object with the dialog.  The
Task object can subsequently be retrieved with @code{gGetTask}.

The value returned is the dialog object passed.
@example
@group
@exdent Example:

object  dlg, task;

gSetTask(dlg, task);
@end group
@end example
@sp 1
See also:  @code{GetTask}
@end deffn




@deffn {SetTitle} SetTitle::Dialog
@sp 2
@example
@group
r = gSetTitle(dlg, title);

object  dlg;     /*  a dialog object    */
char    *title;  /*  new title string   */
object  r;       /*  the dialog obj     */
@end group
@end example
This method is used to set or change the text displayed in the dialog's
title bar.

The value returned is the dialog object passed.
@example
@group
@exdent Example:

object  dlg;

gSetTitle(dlg, "Customer Entry");
@end group
@end example
@end deffn




@deffn {SetTopic} SetTopic::Dialog
@sp 2
@example
@group
pt = gSetTopic(dlg, tpc);

object  dlg;    /*  dialog object         */
char    *tpc;   /*  help topic            */
char    *pt;    /*  previous help topic   */
@end group
@end example
This method is used to associate help text with dialog @code{dlg}.
The help text is defined using the Windows help system and labeled
with the topic indicated by @code{tpc}.  Then, if the user hits
the F1 key while in the dialog, WDS will automatically bring up
the Windows help system and find the indicated topic.

WDS also supports window and control specific topics.  See the appropriate
sections.

This method returns any previous topic associated with the dialog.
@example
@group
@exdent Example:

object  dlg;

gSetTopic(dlg, "myDialogHelp");
@end group
@end example
@sp 1
See also:  The @code{HelpSystem} class and @code{SetTopic::Control} 
@end deffn















@deffn {SetZOrder} SetZOrder::Dialog
@sp 2
@example
@group
r = gSetZOrder(dlg, m);

object  dlg;   /*  a dialog object  */
HWND    m;     /*  Z-order flag     */
object  r;     /*  the dialog obj   */
@end group
@end example
This method is used to set the Z-order position of the dialog window.
The @code{m} argument specifies the position and can be one of the
standard Windows constants such as @code{HWND_TOP}, @code{HWND_TOPMOST},
@code{HWND_BOTTOM}, or @code{HWND_NOTOPMOST}.

The value returned is the dialog object passed.
@example
@group
@exdent Example:

object  dlg;

gSetZOrder(dlg, HWND_TOPMOST);
@end group
@end example
@sp 1
See also:  @code{MoveToTop}
@end deffn




@deffn {TextBrush} TextBrush::Dialog
@sp 2
@example
@group
r = gTextBrush(dlg, brsh);

object  dlg;    /*  a dialog object   */
object  brsh;   /*  brush object      */
object  r;      /*  dlg               */
@end group
@end example
This method is used to set the brush object which is used for foreground
text which is displayed.  Any previously associated brush object will
be disposed.  This brush object will also be automatically disposed when the
dialog is disposed.

The dialog passed is returned.
@example
@group
@exdent Example:

object  dlg;

gTextBrush(dlg, vNew(SolidBrush, 255, 0, 0));
@end group
@end example
@sp 1
See also:  @code{BackBrush, SetTextBrush::Application}
        and the @code{Brush} classes
@end deffn


@deffn {Update} Update::Dialog
@sp 2
@example
@group
r = gUpdate(dlg);

object  dlg;   /*  a dialog object  */
object  r;     /*  the dialog obj   */
@end group
@end example
This method causes all controls in the dialog to push their current
values to their associated data objects.  This is useful when you need
to transfer control values to the underlying data without closing the
dialog.

The value returned is the dialog object passed.
@example
@group
@exdent Example:

object  dlg;

gUpdate(dlg);
@end group
@end example
@end deffn




@deffn {UseDefault} UseDefault::Dialog
@sp 2
@example
@group
r = gUseDefault(dlg);

object  dlg;   /*  a dialog object  */
object  r;     /*  the dialog obj   */
@end group
@end example
This method is used to reset all controls in the dialog to their
default values.

The value returned is the dialog object passed.
@example
@group
@exdent Example:

object  dlg;

gUseDefault(dlg);
@end group
@end example
@end deffn




@deffn {WaitCursor} WaitCursor::Dialog
@sp 2
@example
@group
r = gWaitCursor(dlg, wait);

object  dlg;    /*  a dialog object  */
int     wait;   /*  1=show, 0=hide   */
object  r;      /*  the dialog obj   */
@end group
@end example
This method is used to show or hide the hourglass (wait) cursor for the
dialog.  A reference counting mechanism is used so that nested calls
work correctly --- each call with @code{wait} set to 1 must be balanced
by a corresponding call with @code{wait} set to 0.

The value returned is the dialog object passed.
@example
@group
@exdent Example:

object  dlg;

gWaitCursor(dlg, 1);
/*  long operation  */
gWaitCursor(dlg, 0);
@end group
@end example
@sp 1
See also:  @code{SetCursor}
@end deffn




@subsection Modal Dialogs
Modal dialogs are created with the methods described in this section.
However, since the @code{ModalDialog} class is a subclass of @code{Dialog},
and the @code{Dialog} class implements all the functionality common to
both the @code{ModalDialog} and @code{ModelessDialog} classes, the
majority of the functionality is inherited from and documented in
the @code{Dialog} class.  Therefore, see the @code{Dialog} class
for documentation on additional functionality available to this class.



@deffn {NewDialog} NewDialog::ModalDialog
@sp 2
@example
@group
dlg = mNewDialog(ModalDialog, id, wnd);

unsigned  id;   /*  dialog id      */
object    wnd;  /*  parent window  */
object    dlg;  /*  dialog object  */
@end group
@end example
@code{mNewDialog} is a macro (defined in @code{dynwin.h}) that expands
to a call to @code{gNewDialog} on the @code{Dialog} base class.

This method is used to create a new modal dialog.  The dialog must have
been previously laid out using the resource editor.  The @code{id}
is a macro generated by the resource editor while the programmer defines
the dialog.  This id uniquely identifies a particular dialog.

@code{wnd} allows the specification of the parent window object associated
with the new dialog.  This is an optional parameter and may be @code{NULL}
if it not desired.  If you do specify a parent window, the dialog will
iconize when the window is iconized and will be disposed if the window
is disposed.

The object returned represents the new dialog created and may be
used to further control the dialog.
@example
@group
@exdent Example:

object  wind, dlg;

dlg = mNewDialog(ModalDialog, MY_DIALOG, wind);
@end group
@end example
@sp 1
See also:  @code{NewDialogStr, NewDialogFromFile, Perform::Dialog,}
@iftex
@hfil @break @hglue .63in
@end iftex
@code{NewDialog::ModelessDialog}
@end deffn










@deffn {NewDialogStr} NewDialogStr::ModalDialog
@sp 2
@example
@group
dlg = gNewDialogStr(ModalDialog, id, wnd);

char      *id;  /*  dialog id      */
object    wnd;  /*  parent window  */
object    dlg;  /*  dialog object  */
@end group
@end example
This method is used to create a new modal dialog.  The dialog must have
been previously laid out using the resource editor.  The @code{id} is
the string name associated with the dialog.  This id uniquely identifies
a particular dialog.

@code{wnd} allows the specification of the parent window object associated
with the new dialog.  This is an optional parameter and may be @code{NULL}
if it not desired.  If you do specify a parent window, the dialog will
iconize when the window is iconized and will be disposed if the window
is disposed.

The object returned represents the new dialog created and may be
used to further control the dialog.
@example
@group
@exdent Example:

object  wind, dlg;

dlg = gNewDialogStr(ModalDialog, "mydialog", wind);
@end group
@end example
@sp 1
See also:  @code{NewDialog, NewDialogFromFile, Perform::Dialog,}
@iftex
@hfil @break @hglue .63in
@end iftex
@code{NewDialogStr::ModelessDialog}
@end deffn




@deffn {NewDialogFromFile} NewDialogFromFile::ModalDialog
@sp 2
@example
@group
dlg = gNewDialogFromFile(ModalDialog, file, res, wnd);

char      *file;  /*  CLD file name  */
unsigned  res;    /*  resource id    */
object    wnd;    /*  parent window  */
object    dlg;    /*  dialog object  */
@end group
@end example
This method is used to create a new modal dialog from a CLD resource
file rather than from a compiled Windows resource.  @code{file} is the
path to the CLD file and @code{res} is the resource identifier within
that file.

@code{wnd} allows the specification of the parent window object associated
with the new dialog.  This is an optional parameter and may be @code{NULL}
if it not desired.

The object returned represents the new dialog created and may be
used to further control the dialog.
@example
@group
@exdent Example:

object  wind, dlg;

dlg = gNewDialogFromFile(ModalDialog, "dialogs.cld", 100, wind);
@end group
@end example
@sp 1
See also:  @code{NewDialog, NewDialogStr, Perform::Dialog}
@end deffn









@subsection Modeless Dialogs
Modeless dialogs are created with the methods described in this section.
However, since the @code{ModelessDialog} class is a subclass of
@code{Dialog}, and the @code{Dialog} class implements all the
functionality common to both the @code{ModalDialog} and
@code{ModelessDialog} classes, the majority of the functionality is
inherited from and documented in the @code{Dialog} class.  Therefore,
see the @code{Dialog} class for documentation on additional
functionality available to this class.




@deffn {NewDialog} NewDialog::ModelessDialog
@sp 2
@example
@group
dlg = mNewDialog(ModelessDialog, id, wnd);

unsigned  id;   /*  dialog id      */
object    wnd;  /*  parent window  */
object    dlg;  /*  dialog object  */
@end group
@end example
@code{mNewDialog} is a macro (defined in @code{dynwin.h}) that expands
to a call to @code{gNewDialog} on the @code{Dialog} base class.

This method is used to create a new modeless dialog.  The dialog must have
been previously laid out using the resource editor.  The @code{id}
is a macro generated by the resource editor while the programmer defines
the dialog.  This id uniquely identifies a particular dialog.

@code{wnd} allows the specification of the parent window object associated
with the new dialog.  This is an optional parameter and may be @code{NULL}
if it not desired.  If you do specify a parent window, the dialog will
iconize when the window is iconized and will be disposed if the window
is disposed.

When a new modeless dialog is created, its auto dispose mode is enabled.
This means that, unless disabled, the object representing the dialog
will be automatically disposed when the user closes the window.  This
is done due to the independent nature of modeless dialogs and the need
to be sure that objects no longer needed are disposed.  Note that the
programmer may specify a function which gets executed prior to the
disposal of a modeless dialog in order to capture any data.

The object returned represents the new dialog created and may be
used to further control the dialog.
@example
@group
@exdent Example:

object  wind, dlg;

dlg = mNewDialog(ModelessDialog, MY_DIALOG, wind);
@end group
@end example
@sp 1
See also:  @code{NewDialogStr, NewDialogFromFile, Perform::Dialog,}
@iftex
@hfil @break @hglue .63in
@end iftex
@code{NewDialog::ModalDialog, CompletionFunction::Dialog,}
@iftex
@hfil @break @hglue .63in
@end iftex
@code{AutoDispose::Dialog}
@end deffn












@deffn {NewDialogStr} NewDialogStr::ModelessDialog
@sp 2
@example
@group
dlg = gNewDialogStr(ModelessDialog, id, wnd);

char      *id;  /*  dialog id      */
object    wnd;  /*  parent window  */
object    dlg;  /*  dialog object  */
@end group
@end example
This method is used to create a new modeless dialog.  The dialog must have
been previously laid out using the resource editor.  The @code{id} is
the string name associated with the dialog.  This id uniquely identifies
a particular dialog.

@code{wnd} allows the specification of the parent window object associated
with the new dialog.  This is an optional parameter and may be @code{NULL}
if it not desired.  If you do specify a parent window, the dialog will
iconize when the window is iconized and will be disposed if the window
is disposed.

When a new modeless dialog is created, its auto dispose mode is enabled.
This means that, unless disabled, the object representing the dialog
will be automatically disposed when the user closes the window.  This
is done due to the independent nature of modeless dialogs and the need
to be sure that objects no longer needed are disposed.  Note that the
programmer may specify a function which gets executed prior to the
disposal of a modeless dialog in order to capture any data.

The object returned represents the new dialog created and may be
used to further control the dialog.
@example
@group
@exdent Example:

object  wind, dlg;

dlg = gNewDialogStr(ModelessDialog, "mydialog", wind);
@end group
@end example
@sp 1
See also:  @code{NewDialog, NewDialogFromFile, Perform::Dialog,}
@iftex
@hfil @break @hglue .63in
@end iftex
@code{NewDialogStr::ModalDialog, CompletionFunction::Dialog,}
@iftex
@hfil @break @hglue .63in
@end iftex
@code{AutoDispose::Dialog}
@end deffn




@deffn {NewDialogFromFile} NewDialogFromFile::ModelessDialog
@sp 2
@example
@group
dlg = gNewDialogFromFile(ModelessDialog, file, res, wnd);

char      *file;  /*  CLD file name  */
unsigned  res;    /*  resource id    */
object    wnd;    /*  parent window  */
object    dlg;    /*  dialog object  */
@end group
@end example
This method is used to create a new modeless dialog from a CLD resource
file rather than from a compiled Windows resource.  @code{file} is the
path to the CLD file and @code{res} is the resource identifier within
that file.

@code{wnd} allows the specification of the parent window object associated
with the new dialog.  This is an optional parameter and may be @code{NULL}
if it not desired.

When a new modeless dialog is created, its auto dispose mode is enabled.
This means that, unless disabled, the object representing the dialog
will be automatically disposed when the user closes the window.

The object returned represents the new dialog created and may be
used to further control the dialog.
@example
@group
@exdent Example:

object  wind, dlg;

dlg = gNewDialogFromFile(ModelessDialog, "dialogs.cld", 100, wind);
@end group
@end example
@sp 1
See also:  @code{NewDialog, NewDialogStr, Perform::Dialog,}
@iftex
@hfil @break @hglue .63in
@end iftex
@code{CompletionFunction::Dialog, AutoDispose::Dialog}
@end deffn






@section Controls
The @code{Control} class is never used directly.  It is used to group
functionality common to all the control types which are all subclasses
of this class.  Therefore, this section documents all the methods which
are accessible to all subclasses of this class.


Control instances are normally created via @code{AddControl::Dialog}
and may be subsequently accessed via @code{GetControl::Dialog}.



@subsection  Standard Control Method Arguments

Since all control instance methods have the control object as their
first argument (referred to as @code{ctl}), this argument will not be
described each time.  It always refers to the control object which you
wish to perform the desired operation on.  Each control object will be
an instance of one of the subclasses of the @code{Control} class (such
as @code{TextControl} or @code{NumericControl}).


Many of the methods associated with this class take an argument
identified as @code{id}.  This is a macro defined by the programmer via
the resource editor when the dialog is being defined and is used to
uniquely identify a particular control within the dialog.  Due to the
fact that this argument has the same meaning for every method which
uses it, it will not be defined each time.


@subsection Control Methods












@deffn {AddHandlerAfter} AddHandlerAfter::Control
@sp 2
@example
@group
r = gAddHandlerAfter(ctl, msg, func);

object   ctl;      /*  a control object  */
unsigned msg;      /*  message           */
long    (*func)(); /*  function pointer  */
object  r;         /*  the control obj   */
@end group
@end example
This method is inherited from the @code{Window} class.

This method is used to associate function @code{func} with Windows control
message @code{msg} for control @code{ctl}.  Whenever control @code{ctl}
receives message @code{msg}, @code{func} will be called.

@code{ctl} is the control object who's messages you wish to process.
@code{msg} is the particular message you wish to trap.  These messages
are fully documented in the Windows documentation in the Messages
section.  They normally begin with @code{WM_}.

@code{func} is the function which gets called whenever the specified
message gets received and takes the following form:
@example
@group
long    func(object     ctl,
             HWND       hwnd, 
             UINT       mMsg, 
             WPARAM     wParam, 
             LPARAM     lParam)
@{
        .
        .
        .
        return 0L;  /* or whatever is appropriate  */
@}
@end group
@end example
Where @code{ctl} is the control being sent the message.  The remaining
arguments and return value is fully documented in the Windows documentation
under the @code{WindowProc} function and the Windows Messages documentation.

WDS keeps a list of functions associated with each message associated
with each control.  When a particular message is received the appropriate
list of handler functions gets executed sequentially.
@code{AddHandlerAfter} appends the new function to the end of this list,
and @code{AddHandlerBefore} adds the new function to the beginning of
the list.

WDS may also, and optionally, execute the Windows default procedure
associated with a given message either before or after the user added
list of functions.  This behavior may be controlled via
@code{DefaultProcessingMode}.

Windows will only see the return value of the last message handler executed
including, if applicable, the default.
@example
@group
@exdent Example:

int     hSize, vSize;

static  long    process_wm_size(object  ctl, 
                                HWND    hwnd, 
                                UINT    mMsg, 
                                WPARAM  wParam, 
                                LPARAM  lParam)
@{
        hSize = LOWORD(lParam);
        vSize = HIWORD(lParam);
        return 0L;
@}

        .
        .
        gAddHandlerAfter(ctl, (unsigned) WM_SIZE, process_wm_size);
        .
        .
@end group
@end example
@sp 1
See also:  @code{DefaultProcessingMode, AddHandlerBefore, CallDefaultProc}
@end deffn






@deffn {AddHandlerBefore} AddHandlerBefore::Control
@sp 2
@example
@group
r = gAddHandlerBefore(ctl, msg, func);

object   ctl;      /*  a control object  */
unsigned msg;      /*  message           */
long    (*func)(); /*  function pointer  */
object  r;         /*  the control obj   */
@end group
@end example
This method is inherited from the @code{Window} class.
It is fully documented under @code{AddHandlerAfter}.
@c @example
@c @group
@c @exdent Example:
@c @end group
@c @end example
@sp 1
See also:  @code{AddHandlerAfter}
@end deffn









@deffn {AutoAttach} AutoAttach::Control
@sp 2
@example
@group
prev = gAutoAttach(ctl, mode);

object  ctl;    /*  control object  */
int     mode;   /*  auto-attach mode (0 or 1)  */
int     prev;   /*  previous mode   */
@end group
@end example
This method sets whether the control automatically attaches to
associated database fields.  A non-zero @code{mode} enables
auto-attach; zero disables it.

The value returned is the previous auto-attach mode.
@end deffn



@deffn {CallDefaultProc} CallDefaultProc::Control
@sp 2
@example
@group
r = gCallDefaultProc(ctl, msg, wp, lp)

object   ctl;   /*  a control object  */
unsigned msg;   /*  message           */
WPARAM   wp;    /*  wParam value      */
LPARAM   lp;    /*  lParam value      */
long     r;     /*  result of call    */
@end group
@end example
This method is inherited from the @code{Window} class.

This method is used to explicitly call the default Windows message handler
associated with message @code{msg}.  This method is normally called from
within a programmer defined message handler (see @code{AddHandlerAfter})
which is provided with the @code{wp} and @code{lp} parameters.  These
parameters are fully documented in the Windows Messages documentation
and relate directly to the particular message.

@code{msg} is the particular message you wish to affect.  These messages
are fully documented in the Windows documentation in the Messages
section.  They normally begin with @code{WM_}.
@example
@group
@exdent Example:

gCallDefaultProc(ctl, msg, wParam, lParam);
@end group
@end example
@sp 1
See also:  @code{DefaultProcessingMode, AddHandlerAfter}
@end deffn















@deffn {CheckFunction} CheckFunction::Control
@sp 2
@example
@group
r = gCheckFunction(ctl, fun);

object  ctl;      /*  control object  */
int     (*fun)(); /*  check function  */
object  r;        /*  ctl object      */
@end group
@end example
This method is implemented by all control subclasses (not in the
@code{Control} base class itself).

This method is used to associate an auxiliary checking function to a control.
When a dialog is accepted, WDS goes through each control to assure the
validity of the data associated with each control as defined by the control
object.  The ability to associate an additional function for validating
the data associated with a control gives the programmer the ability to perform
any additional application specific checking against a control's data.

The auxiliary checking function takes the following form:
@example
@group
int     fun(object ctl, object val, char *buf)
@{
        ....
@}
@end group
@end example
Where @code{ctl} is the control object, and @code{val} is a Dynace object which
represents the value the control has associated with it.  @code{buf} is a
pointer to a buffer area which @code{fun} must set to an appropriate error
message if an error occurs.  @code{fun} returns a 1 if an error occurs and
0 otherwise.

All controls support this feature except the @code{PushButton} class,
because it is an immediate action control with no associated value.
@example
@group
@exdent Example:

object  ctl;

gCheckFunction(ctl, fun);
@end group
@end example
@sp 1
See also:  @code{GetControl::Dialog}
@end deffn















@deffn {CheckValue} CheckValue::Control
@sp 2
@example
@group
r = gCheckValue(ctl)

object   ctl;   /*  a control object  */
int      r;     /*  result of call    */
@end group
@end example
This method is used to explicitly cause all error checking functions and
parameters associated with control @code{ctl} to be checked.  If an error
is encountered an appropriate error message will be displayed and a 1
will be returned.  If no error is encountered a 0 is returned.

This method is implemented by all subclasses of @code{Control} and mainly
used internally when the dialog is accepted.
@example
@group
@exdent Example:

gCheckValue(ctl);
@end group
@end example
@sp 1
See also:  @code{CheckFunction::Control}
@end deffn


















@deffn {DeepDispose} DeepDispose::Control
@sp 2
This method performs the same function as @code{Dispose}.  See that
method for details.
@end deffn









@deffn {DefaultProcessingMode} DefaultProcessingMode::Control
@sp 2
@example
@group
r = gDefaultProcessingMode(ctl, msg, mode);

object   ctl;   /*  a control object         */
unsigned msg;   /*  message                  */
int      mode;  /*  default processing mode  */
object   r;     /*  the control obj          */
@end group
@end example
This method is inherited from the @code{Window} class.

This method is used to determine when or if the Windows default message
procedure is processed for a given message (@code{msg}) associated with
a particular control (@code{ctl}).

WDS allows a programmer to specify an arbitrary number of functions
to be executed whenever a control receives a specific message (via
@code{AddHandlerAfter} and @code{AddHandlerBefore}).  Windows has
default procedures associated with many control messages.  At times
it is necessary to replace or augment this default functionality.
@code{DefaultProcessingMode} gives the programmer control over when and
if this default Windows functionality.  @code{mode} is used to specify
the desired mode.  The following table indicates the valid modes:

@table @code
@item 0
Do not execute the Windows default processing
@item 1
Execute default processing @emph{after} programmer defined handlers
@item 2
Execute default processing @emph{before} programmer defined handlers
@end table

Note that the default mode is always @code{1}, and must be explicitly
changed, if desired, for each message associated with each control.

@code{msg} is the particular message you wish to affect.  These messages
are fully documented in the Windows documentation in the Messages
section.  They normally begin with @code{WM_}.
@example
@group
@exdent Example:

gDefaultProcessingMode(ctl, (unsigned) WM_SIZE, 0);
@end group
@end example
@sp 1
See also:  @code{CallDefaultProc, AddHandlerAfter}
@end deffn











@deffn {Disable} Disable::Control
@sp 2
@example
@group
r = gDisable(ctl);

object  ctl;   /*  control object  */
object  r;     /*  ctl             */
@end group
@end example
This method disables the control, graying it out so the user cannot
interact with it.

The value returned is the control object passed.
@end deffn



@deffn {Display} Display::Control
@sp 2
@example
@group
r = gDisplay(ctl);

object  ctl;   /*  control object  */
object  r;     /*  ctl             */
@end group
@end example
This method shows a previously hidden control, making it visible again.

The value returned is the control object passed.
@sp 1
See also:  @code{Hide}
@end deffn



@deffn {Dispose} Dispose::Control
@sp 2
@example
@group
r = gDispose(ctl);

object  ctl;   /*  a control object  */
object  r;     /*  NULL              */
@end group
@end example
This method is used to remove and dispose of a control object when it
is no longer needed.  This method is rarely needed due to the fact that
when a dialog is disposed it automatically calls this method on each of
its associated controls.

The value returned is always @code{NULL} and may be used to null out
the variable which contained the object being disposed in order to
avoid future accidental use.
@example
@group
@exdent Example:

object  ctl;

ctl = gDispose(ctl);
@end group
@end example
@sp 1
See also:  @code{DeepDispose}
@end deffn














@deffn {Enable} Enable::Control
@sp 2
@example
@group
r = gEnable(ctl);

object  ctl;   /*  control object  */
object  r;     /*  ctl             */
@end group
@end example
This method enables the control, making it active and available for user
interaction.  The control may remain disabled if security restrictions
are in effect.

The value returned is the control object passed.
@end deffn



@deffn {GetFont} GetFont::Control
@sp 2
@example
@group
font = gGetFont(ctl);

object  ctl;    /*  control object  */
object  font;   /*  Font object or NULL  */
@end group
@end example
This method returns the @code{Font} object associated with the control,
or @code{NULL} if no font has been set.
@sp 1
See also:  @code{SetFont}
@end deffn



@deffn {GetResourceID} GetResourceID::Control
@sp 2
@example
@group
id = gGetResourceID(ctl);

object    ctl;   /*  control object    */
unsigned  id;    /*  resource identifier  */
@end group
@end example
This method returns the resource identifier associated with the control.
@end deffn



@deffn {GetTabOrder} GetTabOrder::Control
@sp 2
@example
@group
t = gGetTabOrder(ctl);

object  ctl;   /*  control object    */
int     t;     /*  tab order index   */
@end group
@end example
This method returns the current tab order index for the control.
@sp 1
See also:  @code{SetTabOrder}
@end deffn











@deffn {Handle} Handle::Control
@sp 2
@example
@group
h = gHandle(ctl);

object  ctl;    /*  control object   */
HANDLE  h;      /*  Windows handle  */
@end group
@end example
This method is inherited from the @code{Window} class.

This method is used to obtain the Windows internal handle associated with
a control object.  It will only return a valid handle while the dialog
is being performed via @code{gPerform}.  @code{NULL} will be returned
otherwise.

Note that this method may be used with most WDS objects in order to obtain
the internal handle that Windows normally associates with each type of object.
See the appropriate documentation.
@example
@group
@exdent Example:

object  ctl;
HANDLE  h;

h = gHandle(ctl);
@end group
@end example
@sp 1
See also:  @code{GetControl::Dialog}
@end deffn











@deffn {Hide} Hide::Control
@sp 2
@example
@group
r = gHide(ctl);

object  ctl;   /*  control object  */
object  r;     /*  ctl             */
@end group
@end example
This method hides the control so that it is not visible on the dialog.

The value returned is the control object passed.
@sp 1
See also:  @code{Display}
@end deffn



@deffn {SetFocus} SetFocus::Control
@sp 2
@example
@group
r = gSetFocus(ctl);

object  ctl;   /*  control object  */
object  r;     /*  ctl             */
@end group
@end example
This method sets the keyboard focus to the specified control.

The value returned is the control object passed.
@end deffn



@deffn {SetFont} SetFont::Control
@sp 2
@example
@group
r = gSetFont(ctl, font);

object  ctl;    /*  control object  */
object  font;   /*  Font object     */
object  r;      /*  ctl             */
@end group
@end example
This method sets the display font used by the control.

The value returned is the control object passed.
@sp 1
See also:  @code{GetFont}
@end deffn



@deffn {SetName} SetName::Control
@sp 2
@example
@group
r = gSetName(ctl, name);

object  ctl;    /*  control object   */
char    *name;  /*  control name     */
object  r;      /*  ctl              */
@end group
@end example
This method sets the name of the control.  The name is used for
programmatic lookup of controls within a dialog.

The value returned is the control object passed.
@end deffn



@deffn {SetTabOrder} SetTabOrder::Control
@sp 2
@example
@group
prev = gSetTabOrder(ctl, t);

object  ctl;    /*  control object     */
int     t;      /*  tab order index    */
int     prev;   /*  previous index     */
@end group
@end example
This method sets the tab order index for the control.  The tab order
determines the sequence in which controls receive focus when the user
presses the Tab key.

The value returned is the previous tab order index.
@sp 1
See also:  @code{GetTabOrder}
@end deffn



@deffn {SetTopic} SetTopic::Control
@sp 2
@example
@group
pt = gSetTopic(ctl, tpc);

object  ctl;    /*  control object        */
char    *tpc;   /*  help topic            */
char    *pt;    /*  previous help topic   */
@end group
@end example
This method is inherited from the @code{Window} class.

This method is used to associate help text with specific control @code{ctl}.
The help text is defined using the Windows help system and labeled
with the topic indicated by @code{tpc}.  Then, if the user hits
the F1 key while in the control, WDS will automatically bring up
the Windows help system and find the indicated topic.  If no control
specific help is specified, WDS will use any help associated with the
dialog as a whole.

WDS also supports dialog and control specific topics.  See the appropriate
sections.

This method returns any previous topic associated with the control.
@example
@group
@exdent Example:

object  ctl;

gSetTopic(ctl, "myControlHelp");
@end group
@end example
@sp 1
See also:  The @code{HelpSystem} class and @code{SetTopic::Dialog}
@end deffn



@deffn {Update} Update::Control
@sp 2
@example
@group
r = gUpdate(ctl);

object  ctl;   /*  control object  */
object  r;     /*  ctl             */
@end group
@end example
This method pushes the control's current value to its associated
data object.

The value returned is the control object passed.
@end deffn



@deffn {UseDefault} UseDefault::Control
@sp 2
@example
@group
r = gUseDefault(ctl);

object  ctl;   /*  control object  */
object  r;     /*  ctl             */
@end group
@end example
This method resets the control to its default value.

The value returned is the control object passed.
@end deffn



@subsection Text Control
The @code{TextControl} class is a control which enables the user to
enter arbitrary alphanumeric data on a single line.  This control
has numerous facilities to control the length and type of input the
user is allowed to enter.

The @code{TextControl} class is a subclass of @code{Control} and as such
inherits all the functionality associated with that class.  This section
only documents functionality particular to this class.

Standard control arguments are documented in the section entitled
``Standard Control Method Arguments''.












@deffn {Attach} Attach::TextControl
@sp 2
@example
@group
r = gAttach(ctl, val);

object  ctl;   /*  control object  */
object  val;   /*  ctl value       */
object  r;     /*  control object  */
@end group
@end example
This method is used to associate an independent Dynace object with the
value associated with a control object.  @code{val} should be an
instance of the @code{String} class and will be automatically updated
to reflect the value associated with the control.

This object (@code{val}) will never be disposed by WDS, even after
the control or associated dialog are disposed.  Therefore, this
is one way of gaining access to a control's value after the life
of the control.  It is the programmer's responsibility to dispose of
the object when it is no longer needed.
@example
@group
@exdent Example:

object  ctl, val;

val = gNew(String);
gAttach(ctl, val);
@end group
@end example
@sp 1
See also:  @code{Value, StringValue, SetValue}
@end deffn












@deffn {Capitalize} Capitalize::TextControl
@sp 2
@example
@group
r = gCapitalize(ctl);

object  ctl;   /*  control object  */
object  r;     /*  control object  */
@end group
@end example
This method is used to enable the auto-capitalize feature of the text
control.  If enabled, all alpha entry made by the user will be converted
to upper case both on the display and internally to the control.
@example
@group
@exdent Example:

object  ctl;

gCapitalize(ctl);
@end group
@end example
@sp 1
See also:  @code{TextRange}
@end deffn



@deffn {GetMinLength} GetMinLength::TextControl
@sp 2
@example
@group
len = gGetMinLength(ctl);

object  ctl;   /*  control object        */
int     len;   /*  minimum input length  */
@end group
@end example
This method returns the minimum required input length for the control.
@sp 1
See also:  @code{MinLength}
@end deffn



@deffn {HideData} HideData::TextControl
@sp 2
@example
@group
prev = gHideData(ctl, flag);

object  ctl;    /*  control object  */
int     flag;   /*  hide data flag  */
int     prev;   /*  previous flag   */
@end group
@end example
This method sets whether the data displayed in the control is hidden.
A non-zero @code{flag} hides the data; zero shows it.

The value returned is the previous hide data flag.
@end deffn













@deffn {MaxLength} MaxLength::TextControl
@sp 2
@example
@group
r = gMaxLength(ctl, len);

object  ctl;   /*  control object  */
int     len;   /*  maximum length  */
object  r;     /*  control object  */
@end group
@end example
This method is used to set the maximum number of characters the user may
enter into a text control.  If the user attempts to enter more than
@code{len} characters, including spaces, the system will just beep.  If
the number of characters exceeds the field display width, the display
will scroll.
@example
@group
@exdent Example:

object  ctl;

gMaxLength(ctl, 25);
@end group
@end example
@sp 1
See also:  @code{MinLength, TextRange, Capitalize, CheckFunction::Control}
@end deffn












@deffn {MinLength} MinLength::TextControl
@sp 2
@example
@group
r = gMinLength(ctl, len);

object  ctl;   /*  control object  */
int     len;   /*  minimum length  */
object  r;     /*  control object  */
@end group
@end example
This method is used to set the minimum number of characters the user must
enter into a text control.  If the user attempts to accept the dialog
without entering @code{len} characters, including spaces, WDS will
issue an error message and return the user to the incomplete field.
@example
@group
@exdent Example:

object  ctl;

gMinLength(ctl, 5);
@end group
@end example
@sp 1
See also:  @code{MaxLength, TextRange, Capitalize, CheckFunction::Control}
@end deffn









@deffn {NewCtl} NewCtl::TextControl
@sp 2
@example
@group
ctl = mNewCtl(TextControl, id);

unsigned  id;   /*  control id      */
object   ctl;   /*  control object  */
@end group
@end example
@code{mNewCtl} is a macro (defined in @code{dynwin.h}) that expands to
a call to @code{vNew}.

This method is used to create a new control object to be identified as
@code{id} (see section ``Standard Control Method Arguments'').  This
method is mainly used internally.  A programmer would more often
use @code{AddControl::Dialog} to create controls and associate them
to a dialog.
@example
@group
@exdent Example:

object  ctl;

ctl = mNewCtl(TextControl, PERSON_NAME);
@end group
@end example
@sp 1
See also:  @code{AddControl::Dialog, GetControl::Dialog}
@end deffn



@deffn {PasswordFlag} PasswordFlag::TextControl
@sp 2
@example
@group
flag = gPasswordFlag(ctl);

object  ctl;    /*  control object      */
int     flag;   /*  current password flag  */
@end group
@end example
This method returns the current password flag for the control.
@sp 1
See also:  @code{SetPasswordFlag}
@end deffn



@deffn {Select} Select::TextControl
@sp 2
@example
@group
r = gSelect(ctl);

object  ctl;   /*  control object  */
object  r;     /*  ctl             */
@end group
@end example
This method selects (highlights) all text in the control.

The value returned is the control object passed.
@end deffn



@deffn {SetMask} SetMask::TextControl
@sp 2
@example
@group
r = gSetMask(ctl, mask);

object  ctl;    /*  control object  */
char    *mask;  /*  input mask pattern or NULL  */
object  r;      /*  ctl             */
@end group
@end example
This method sets an input mask pattern for the control.  The mask
defines the format of allowable input.  For example,
@code{"(###) ###-####"} would be used for a phone number field, where
@code{#} represents a digit.  Passing @code{NULL} clears the mask.

The value returned is the control object passed.
@end deffn



@deffn {SetPasswordFlag} SetPasswordFlag::TextControl
@sp 2
@example
@group
prev = gSetPasswordFlag(ctl, flag);

object  ctl;    /*  control object     */
int     flag;   /*  password mode flag */
int     prev;   /*  previous flag      */
@end group
@end example
This method sets password entry mode for the control.  When enabled
(non-zero @code{flag}), characters typed by the user are displayed as
asterisks.

The value returned is the previous password flag.
@sp 1
See also:  @code{PasswordFlag}
@end deffn










@deffn {SetStringValue} SetStringValue::TextControl
@sp 2
@example
@group
r = gSetStringValue(ctl, val);

object  ctl;    /*  control object  */
char    *val;   /*  ctl value       */
object  r;      /*  control object  */
@end group
@end example
This method is used to set the value associated with a control.  It is
often used to set the initial value prior to performing a dialog.
@code{val} should be a pointer to a character string which represents
the desired value for the control.

Any previously associated value object will be disposed when the new
value is set.  
@example
@group
@exdent Example:

object  ctl;

gSetStringValue(ctl, "Some value");
@end group
@end example
@sp 1
See also:  @code{SetValue, StringValue}
@end deffn








@deffn {SetValue} SetValue::TextControl
@sp 2
@example
@group
r = gSetValue(ctl, val);

object  ctl;    /*  control object  */
object  val;    /*  ctl value       */
object  r;      /*  control object  */
@end group
@end example
This method is used to set the value associated with a control.  It is
often used to set the initial value prior to performing a dialog.
@code{val} should be a Dynace object which is an instance of the
@code{String} class and initialized to the value desired for the control.

Any previously associated object will be disposed when the new value is set.
Also, @code{val} will automatically be disposed when the control or associated
dialog is disposed.
@example
@group
@exdent Example:

object  ctl;

gSetValue(ctl, gNewWithStr(String, "Some value"));
@end group
@end example
@sp 1
See also:  @code{SetStringValue, Value, Attach}
@end deffn









@deffn {StringValue} StringValue::TextControl
@sp 2
@example
@group
val = gStringValue(ctl);

object  ctl;   /*  control object  */
char    *val;  /*  ctl value       */
@end group
@end example
This method is used to obtain a C character pointer which represents the
value associated with the control.  This pointer will not be valid
after the control or associated dialog are disposed.
@example
@group
@exdent Example:

object  ctl;
char    *val;

val = gStringValue(ctl);
@end group
@end example
@sp 1
See also:  @code{Value, Attach, SetStringValue, CtlStringValue::Dialog}
@end deffn














@deffn {TextRange} TextRange::TextControl
@sp 2
@example
@group
r = gTextRange(ctl, min, max);

object  ctl;   /*  control object  */
int     min;   /*  minimum length  */
int     max;   /*  maximum length  */
object  r;     /*  control object  */
@end group
@end example
This method is used to set the minimum number of characters the user must
enter into a text control as well as the maximum number of characters
allowed.

If the user attempts to enter more than @code{max} characters, including
spaces, the system will just beep.  If the number of characters exceeds
the field display width, the display will scroll.

If the user attempts to accept the dialog without entering @code{min}
characters, including spaces, WDS will issue an error message and return
the user to the incomplete field.
@example
@group
@exdent Example:

object  ctl;

gTextRange(ctl, 5, 25);
@end group
@end example
@sp 1
See also:  @code{MaxLength, MinLength, Capitalize, CheckFunction::Control}
@end deffn











@deffn {Value} Value::TextControl
@sp 2
@example
@group
val = gValue(ctl);

object  ctl;   /*  control object  */
object  val;   /*  ctl value       */
@end group
@end example
This method is used to obtain a Dynace object which represents the value
associated with the control.  The object returned will be an instance of
the @code{String} class.  This object will be disposed by WDS when the
control object or associated dialog is disposed.
@example
@group
@exdent Example:

object  ctl, val;

val = gValue(ctl);
@end group
@end example
@sp 1
See also:  @code{StringValue, Attach, SetValue}
@end deffn








@subsection Numeric Control
The @code{NumericControl} class is a control which enables the user to
enter arbitrary numeric data.  This control has numerous facilities to
control the range and type of input the user is allowed to enter.

The @code{NumericControl} class is a subclass of @code{Control} and as such
inherits all the functionality associated with that class.  This section
only documents functionality particular to this class.

Standard control arguments are documented in the section entitled
``Standard Control Method Arguments''.












@deffn {Attach} Attach::NumericControl
@sp 2
@example
@group
r = gAttach(ctl, val);

object  ctl;   /*  control object  */
object  val;   /*  ctl value       */
object  r;     /*  control object  */
@end group
@end example
This method is used to associate an independent Dynace object with the
value associated with a control object.  @code{val} should be an
instance of one of the Dynace numeric classes which are subclasses of
the @code{Number} class, such as @code{ShortInteger} or
@code{DoubleFloat}.  This numeric object will be automatically updated
to reflect the value associated with the control.

This object (@code{val}) will never be disposed by WDS, even after
the control or associated dialog are disposed.  Therefore, this
is one way of gaining access to a control's value after the life
of the control.  It is the programmer's responsibility to dispose of
the object when it is no longer needed.
@example
@group
@exdent Example:

object  ctl, val;

val = gNewWithDouble(DoubleFloat, 0.0);
gAttach(ctl, val);
@end group
@end example
@sp 1
See also:  @code{Value, DoubleValue, ShortValue, SetValue}
@end deffn












@deffn {DoubleValue} DoubleValue::NumericControl
@sp 2
@example
@group
val = gDoubleValue(ctl);

object  ctl;   /*  control object  */
double  val;   /*  ctl value       */
@end group
@end example
This method is used to obtain a C double which represents the
value associated with the control.  
@example
@group
@exdent Example:

object  ctl;
double  val;

val = gDoubleValue(ctl);
@end group
@end example
@sp 1
See also:  @code{Value, Attach, SetDoubleValue, CtlDoubleValue::Dialog}
@iftex
@hfil @break @hglue .63in 
@end iftex
@code{UnsignedShortValue, ShortValue, LongValue}
@end deffn











@deffn {LongValue} LongValue::NumericControl
@sp 2
@example
@group
val = gLongValue(ctl);

object  ctl;   /*  control object  */
long    val;   /*  ctl value       */
@end group
@end example
This method is used to obtain a C long integer which represents the
value associated with the control.  
@example
@group
@exdent Example:

object  ctl;
long    val;

val = gLongValue(ctl);
@end group
@end example
@sp 1
See also:  @code{Value, Attach, SetLongValue, CtlLongValue::Dialog}
@iftex
@hfil @break @hglue .63in 
@end iftex
@code{UnsignedShortValue, ShortValue, DoubleValue}
@end deffn








@deffn {NewCtl} NewCtl::NumericControl
@sp 2
@example
@group
ctl = mNewCtl(NumericControl, id);

unsigned  id;   /*  control id      */
object   ctl;   /*  control object  */
@end group
@end example
@code{mNewCtl} is a macro (defined in @code{dynwin.h}) that expands to
a call to @code{vNew}.

This method is used to create a new control object to be identified as
@code{id} (see section ``Standard Control Method Arguments'').  This
method is mainly used internally.  A programmer would more often
use @code{AddControl::Dialog} to create controls and associate them
to a dialog.
@example
@group
@exdent Example:

object  ctl;

ctl = mNewCtl(NumericControl, PERSON_AGE);
@end group
@end example
@sp 1
See also:  @code{AddControl::Dialog, GetControl::Dialog}
@end deffn









@deffn {NumericRange} NumericRange::NumericControl
@sp 2
@example
@group
r = gNumericRange(ctl, min, max, dp);

object  ctl;   /*  control object  */
double  min;   /*  minimum value   */
double  max;   /*  maximum value   */
int     dp;    /*  decimal places  */
object  r;     /*  control object  */
@end group
@end example
This method is used to set the minimum and maximum allowable values
the control will accept from the user.  It also sets the number of
digits to the right of the decimal point the control will accept.

If the user enters a number outside the specified range, WDS will
issue an error message and prompt the user for an acceptable value.
@example
@group
@exdent Example:

object  ctl;

gNumericRange(ctl, 5.0, 25.0, 1);
@end group
@end example
@sp 1
See also:  @code{CheckFunction::Control}
@end deffn



@deffn {Select} Select::NumericControl
@sp 2
@example
@group
r = gSelect(ctl);

object  ctl;   /*  control object  */
object  r;     /*  ctl             */
@end group
@end example
This method selects (highlights) all text in the numeric control.

The value returned is the control object passed.
@end deffn



@deffn {SetDefaultDouble} SetDefaultDouble::NumericControl
@sp 2
@example
@group
r = gSetDefaultDouble(ctl, val);

object  ctl;    /*  control object   */
double  val;    /*  default value    */
object  r;      /*  ctl              */
@end group
@end example
This method sets the default double value for the numeric control.
This value is used when the control is reset via @code{UseDefault}.

The value returned is the control object passed.
@sp 1
See also:  @code{UseDefault, SetDefaultLong}
@end deffn



@deffn {SetDefaultLong} SetDefaultLong::NumericControl
@sp 2
@example
@group
r = gSetDefaultLong(ctl, val);

object  ctl;   /*  control object   */
long    val;   /*  default value    */
object  r;     /*  ctl              */
@end group
@end example
This method sets the default long integer value for the numeric control.
This value is used when the control is reset via @code{UseDefault}.

The value returned is the control object passed.
@sp 1
See also:  @code{UseDefault, SetDefaultDouble}
@end deffn














@deffn {SetDoubleValue} SetDoubleValue::NumericControl
@sp 2
@example
@group
r = gSetDoubleValue(ctl, val);

object  ctl;    /*  control object  */
double  val;    /*  ctl value       */
object  r;      /*  control object  */
@end group
@end example
This method is used to set the value associated with a control.  It is
often used to set the initial value prior to performing a dialog.
@code{val} should be the value which represents the desired value for
the control.

Any previously associated value object (associated via @code{SetValue})
will be disposed when the new value is set.
@example
@group
@exdent Example:

object  ctl;

gSetDoubleValue(ctl, 3.14159265358979);
@end group
@end example
@sp 1
See also:  @code{SetValue, DoubleValue, CtlDoubleValue::Dialog, SetUShortValue,}
@iftex
@hfil @break @hglue .64in   
@end iftex
@code{SetShortValue, SetLongValue}
@end deffn



@deffn {SetFormat} SetFormat::NumericControl
@sp 2
@example
@group
r = gSetFormat(ctl, fmt, wth);

object  ctl;    /*  control object          */
char    *fmt;   /*  printf-style format     */
int     wth;    /*  display width           */
object  r;      /*  ctl                     */
@end group
@end example
This method sets the display format and width for the numeric control.
@code{fmt} is a @code{printf}-style format string (e.g., @code{"%%.2f"})
and @code{wth} specifies the display width.

The value returned is the control object passed.
@end deffn















@deffn {SetLongValue} SetLongValue::NumericControl
@sp 2
@example
@group
r = gSetLongValue(ctl, val);

object  ctl;    /*  control object  */
long    val;    /*  ctl value       */
object  r;      /*  control object  */
@end group
@end example
This method is used to set the value associated with a control.  It is
often used to set the initial value prior to performing a dialog.
@code{val} should be the value which represents the desired value for
the control.

Any previously associated value object (associated via @code{SetValue})
will be disposed when the new value is set.
@example
@group
@exdent Example:

object  ctl;

gSetLongValue(ctl, 66L);
@end group
@end example
@sp 1
See also:  @code{SetValue, LongValue, CtlLongValue::Dialog, SetUShortValue,}
@iftex
@hfil @break @hglue .63in 
@end iftex
@code{SetShortValue, SetDoubleValue}
@end deffn














@deffn {SetShortValue} SetShortValue::NumericControl
@sp 2
@example
@group
r = gSetShortValue(ctl, val);

object  ctl;    /*  control object  */
int     val;    /*  ctl value       */
object  r;      /*  control object  */
@end group
@end example
This method is used to set the value associated with a control.  It is
often used to set the initial value prior to performing a dialog.
@code{val} should be the value which represents the desired value for
the control.

Any previously associated value object (associated via @code{SetValue})
will be disposed when the new value is set.
@example
@group
@exdent Example:

object  ctl;

gSetShortValue(ctl, 66);
@end group
@end example
@sp 1
See also:  @code{SetValue, ShortValue, CtlShortValue::Dialog, SetUShortValue,}
@iftex
@hfil @break @hglue .63in 
@end iftex
@code{SetLongValue, SetDoubleValue}
@end deffn











@deffn {SetUShortValue} SetUShortValue::NumericControl
@sp 2
@example
@group
r = gSetUShortValue(ctl, val);

object    ctl;  /*  control object  */
unsigned  val;  /*  ctl value       */
object    r;    /*  control object  */
@end group
@end example
This method is used to set the value associated with a control.  It is
often used to set the initial value prior to performing a dialog.
@code{val} should be the value which represents the desired value for
the control.

Any previously associated value object (associated via @code{SetValue})
will be disposed when the new value is set.
@example
@group
@exdent Example:

object  ctl;

gSetUShortValue(ctl, 66);
@end group
@end example
@sp 1
See also:  @code{SetValue, UnsignedShortValue, CtlUnsignedShortValue::Dialog}
@iftex
@hfil @break @hglue .63in 
@end iftex
@code{SetShortValue, SetLongValue, SetDoubleValue}
@end deffn
















@deffn {SetValue} SetValue::NumericControl
@sp 2
@example
@group
r = gSetValue(ctl, val);

object  ctl;    /*  control object  */
object  val;    /*  ctl value       */
object  r;      /*  control object  */
@end group
@end example
This method is used to set the value associated with a control.  It is
often used to set the initial value prior to performing a dialog.
@code{val} should be a Dynace object which is an instance of one of the
subclasses of the @code{Number} class and initialized to the value
desired for the control.

Any previously associated object will be disposed when the new value is set.
Also, @code{val} will automatically be disposed when the control or associated
dialog is disposed.
@example
@group
@exdent Example:

object  ctl;

gSetValue(ctl, gNewWithDouble(DoubleFloat, 3.141));
@end group
@end example
@sp 1
See also:  @code{Value, SetShortValue, SetUShortValue}
        @code{SetLongValue, SetDoubleValue}
@end deffn











@deffn {ShortValue} ShortValue::NumericControl
@sp 2
@example
@group
val = gShortValue(ctl);

object  ctl;   /*  control object  */
short   val;   /*  ctl value       */
@end group
@end example
This method is used to obtain a C short integer which represents the
value associated with the control.  
@example
@group
@exdent Example:

object  ctl;
short   val;

val = gShortValue(ctl);
@end group
@end example
@sp 1
See also:  @code{Value, Attach, SetShortValue, CtlShortValue::Dialog}
@iftex
@hfil @break @hglue .63in 
@end iftex
@code{UnsignedShortValue, LongValue, DoubleValue}
@end deffn











@deffn {UnsignedShortValue} UnsignedShortValue::NumericControl
@sp 2
@example
@group
val = gUnsignedShortValue(ctl);

object  ctl;           /*  control object  */
unsigned short  val;   /*  ctl value       */
@end group
@end example
This method is used to obtain a C unsigned short integer which represents the
value associated with the control.  
@example
@group
@exdent Example:

object  ctl;
unsigned short  val;

val = gUnsignedShortValue(ctl);
@end group
@end example
@sp 1
See also:  @code{Value, Attach, SetUShortValue, CtlUnsignedShortValue::Dialog}
@iftex
@hfil @break @hglue .63in 
@end iftex
@code{ShortValue, LongValue, DoubleValue}
@end deffn










@deffn {Value} Value::NumericControl
@sp 2
@example
@group
val = gValue(ctl);

object  ctl;   /*  control object  */
object  val;   /*  ctl value       */
@end group
@end example
This method is used to obtain a Dynace object which represents the value
associated with the control.  The object returned will be an instance of
one of the subclasses of the @code{Number} class.  This object will be
disposed by WDS when the control object or associated dialog is
disposed.
@example
@group
@exdent Example:

object  ctl, val;

val = gValue(ctl);
@end group
@end example
@sp 1
See also:  @code{ShortValue, DoubleValue, Attach, SetValue}
@end deffn









@subsection Date Control
The @code{DateControl} class is a control which enables the user to
enter date information.  The user enters dates in the form mm/dd/yy
or mm/dd/yyyy, and when the user leaves the field it changes to the
form MMM DD, YYYY.  For example, the user would enter 6/8/59 and
when the field is accepted it will be changed to display Jun 8, 1959.

Dates are stored internally as C longs in the form YYYYMMDD so that
6/8/59 would be represented as the long 19590608.  This control has
numerous facilities to control the range and type of input the user is
allowed to enter.

The @code{DateControl} class is a subclass of @code{Control} and as such
inherits all the functionality associated with that class.  This section
only documents functionality particular to this class.

Standard control arguments are documented in the section entitled
``Standard Control Method Arguments''.












@deffn {Attach} Attach::DateControl
@sp 2
@example
@group
r = gAttach(ctl, val);

object  ctl;   /*  control object  */
object  val;   /*  ctl value       */
object  r;     /*  control object  */
@end group
@end example
This method is used to associate an independent Dynace object with the
value associated with a control object.  @code{val} should be an
instance of the @code{Date} class.  This date object will be
automatically updated to reflect the value associated with the control.

This object (@code{val}) will never be disposed by WDS, even after
the control or associated dialog are disposed.  Therefore, this
is one way of gaining access to a control's value after the life
of the control.  It is the programmer's responsibility to dispose of
the object when it is no longer needed.
@example
@group
@exdent Example:

object  ctl, val;

val = vNew(Date, 0L);
gAttach(ctl, val);
@end group
@end example
@sp 1
See also:  @code{Value, LongValue, SetValue, SetLongValue}
@end deffn
















@deffn {DateRange} DateRange::DateControl
@sp 2
@example
@group
r = gDateRange(ctl, min, max, an);

object  ctl;   /*  control object  */
long    min;   /*  minimum value   */
long    max;   /*  maximum value   */
int     an;    /*  allow none      */
object  r;     /*  control object  */
@end group
@end example
This method is used to set the minimum and maximum allowable values
the control will accept from the user.  It also controls whether
or not the control will accept no date.

@code{min} and @code{max} are of the form YYYYMMDD.  For example
6/8/59 would be represented as 19590608.  If @code{an} is 0
the user must enter a valid date between @code{min} and @code{max}.
If, however, @code{an} is 1, the user must either enter a valid
date between the given ranges or may also leave the field blank.
If the user does not enter a date the field value will be 0L.

If the user enters an invalid date or one outside the specified range,
WDS will issue an error message and prompt the user for an acceptable
value.
@example
@group
@exdent Example:

object  ctl;

gDateRange(ctl, 19990101L, 19991231L, 1);
@end group
@end example
@sp 1
See also:  @code{CheckFunction::Control}
@end deffn








@deffn {LongValue} LongValue::DateControl
@sp 2
@example
@group
val = gLongValue(ctl);

object  ctl;   /*  control object  */
long    val;   /*  ctl value       */
@end group
@end example
This method is used to obtain a C long integer which represents the
value associated with the control.  It will be in the form YYYYMMDD.
For example 6/8/59 would be represented as 19590608.
@example
@group
@exdent Example:

object  ctl;
long    val;

val = gLongValue(ctl);
@end group
@end example
@sp 1
See also:  @code{Value, Attach, SetLongValue, CtlLongValue::Dialog}
@end deffn










@deffn {NewCtl} NewCtl::DateControl
@sp 2
@example
@group
ctl = mNewCtl(DateControl, id);

unsigned  id;   /*  control id      */
object   ctl;   /*  control object  */
@end group
@end example
@code{mNewCtl} is a macro (defined in @code{dynwin.h}) that expands to
a call to @code{vNew}.

This method is used to create a new control object to be identified as
@code{id} (see section ``Standard Control Method Arguments'').  This
method is mainly used internally.  A programmer would more often
use @code{AddControl::Dialog} to create controls and associate them
to a dialog.
@example
@group
@exdent Example:

object  ctl;

ctl = mNewCtl(DateControl, PERSON_BD);
@end group
@end example
@sp 1
See also:  @code{AddControl::Dialog, GetControl::Dialog}
@end deffn










@deffn {SetLongValue} SetLongValue::DateControl
@sp 2
@example
@group
r = gSetLongValue(ctl, val);

object  ctl;    /*  control object  */
long    val;    /*  ctl value       */
object  r;      /*  control object  */
@end group
@end example
This method is used to set the value associated with a control.  It is
often used to set the initial value prior to performing a dialog.
@code{val} should be the value which represents the desired value for
the control of the form YYYYMMDD.  For example 6/8/59 would be
19590608.

Any previously associated value object (associated via @code{SetValue})
will be disposed when the new value is set.
@example
@group
@exdent Example:

object  ctl;

gSetLongValue(ctl, 19990101L);
@end group
@end example
@sp 1
See also:  @code{SetValue, LongValue, CtlLongValue::Dialog}
@end deffn













@deffn {SetValue} SetValue::DateControl
@sp 2
@example
@group
r = gSetValue(ctl, val);

object  ctl;    /*  control object  */
object  val;    /*  ctl value       */
object  r;      /*  control object  */
@end group
@end example
This method is used to set the value associated with a control.  It is
often used to set the initial value prior to performing a dialog.
@code{val} should be a Dynace object which is an instance of the
@code{Date} class and initialized to the value desired for the control.

Any previously associated object will be disposed when the new value is set.
Also, @code{val} will automatically be disposed when the control or associated
dialog is disposed.
@example
@group
@exdent Example:

object  ctl;

gSetValue(ctl, vNew(Date, 19940101L));
@end group
@end example
@sp 1
See also:  @code{SetLongValue, Value}
@end deffn














@deffn {Value} Value::DateControl
@sp 2
@example
@group
val = gValue(ctl);

object  ctl;   /*  control object  */
object  val;   /*  ctl value       */
@end group
@end example
This method is used to obtain a Dynace object which represents the value
associated with the control.  The object returned will be an instance of
the @code{Date} class.  This object will be disposed by WDS when the
control object or associated dialog is disposed.
@example
@group
@exdent Example:

object  ctl, val;

val = gValue(ctl);
@end group
@end example
@sp 1
See also:  @code{LongValue, Attach, SetValue}
@end deffn

















@subsection Push Buttons
The @code{PushButton} class represents a control which allows the user
to perform an immediate action.  That means that there is no particular
data associated with the control.  The user pushes the button and an
associated action gets performed at that point.  The @code{PushButton}
class allows the programmer to associate a C function to the control
such that whenever the user clicks on the button the specified function
will be evoked.

The @code{PushButton} class is a subclass of @code{Control} and as such
inherits all the functionality associated with that class.  This section
only documents functionality particular to this class.

Standard control arguments are documented in the section entitled
``Standard Control Method Arguments''.





@deffn {Click} Click::PushButton
@sp 2
@example
@group
r = gClick(ctl);

object  ctl;   /*  push button object  */
object  r;     /*  ctl                 */
@end group
@end example
This method programmatically simulates a button click, causing the
button's associated function to be executed as if the user had
clicked it.

The value returned is the control object passed.
@end deffn



@deffn {GetExitType} GetExitType::PushButton
@sp 2
@example
@group
type = gGetExitType(ctl);

object  ctl;    /*  push button object  */
int     type;   /*  exit type           */
@end group
@end example
This method returns the button's current exit type.
@sp 1
See also:  @code{SetExitType}
@end deffn



@deffn {NewCtl} NewCtl::PushButton
@sp 2
@example
@group
ctl = mNewCtl(PushButton, id);

unsigned  id;   /*  control id      */
object   ctl;   /*  control object  */
@end group
@end example
@code{mNewCtl} is a macro (defined in @code{dynwin.h}) that expands to
a call to @code{vNew}.

This method is used to create a new control object to be identified as
@code{id} (see section ``Standard Control Method Arguments'').  This
method is mainly used internally.  A programmer would more often
use @code{AddControl::Dialog} to create controls and associate them
to a dialog.
@example
@group
@exdent Example:

object  ctl;

ctl = mNewCtl(PushButton, MY_BUTTON);
@end group
@end example
@sp 1
See also:  @code{AddControl::Dialog, GetControl::Dialog}
@end deffn









@deffn {Perform} Perform::PushButton
@sp 2
@example
@group
r = gPerform(ctl);

object  ctl;      /*  control object  */
int     r;        /*  result          */
@end group
@end example
This method is used to manually evoke the function a programmer associated 
with a push button.  The value returned is the result of the function
attached to the button.

This method is seldom needed because WDS automatically evokes the function
when the user clicks the button.
@example
@group
@exdent Example:

object  ctl;

gPerform(ctl);
@end group
@end example
@sp 1
See also:  @code{SetFunction}
@end deffn



@deffn {SetExitType} SetExitType::PushButton
@sp 2
@example
@group
r = gSetExitType(ctl, type);

object  ctl;    /*  push button object  */
int     type;   /*  exit type           */
object  r;      /*  ctl                 */
@end group
@end example
This method sets the exit type for the button, which determines the
dialog close behavior when the button is pressed.

The value returned is the control object passed.
@sp 1
See also:  @code{GetExitType}
@end deffn






@deffn {SetFunction} SetFunction::PushButton
@sp 2
@example
@group
r = gSetFunction(ctl, fun);

object  ctl;      /*  control object           */
int     (*fun)(); /*  function pointer         */
ofun    r;        /*  previous function        */
@end group
@end example
This method is inherited from the @code{Control} class.

This method is used to associate a C function to a push button such that
if the user clicks the button the C function will immediately get
evoked.  The value returned is the previous function, if any, associated
with the control.

The C function takes the following form:
@example
@group
int     fun(object ctl, object dlg)
@{
        ....
@}
@end group
@end example
Where @code{ctl} is the control object, and @code{dlg} is the object
representing the dialog the control is associated with.  The value
returned by @code{fun} only has meaning if the button was either
@code{IDOK} or @code{IDCANCEL} (the ones normally used to exit the
dialog).  If the return value is 1 then the dialog will exit like normal.
If the return value is 0 then the dialog will not exit.
@example
@group
@exdent Example:

object  ctl;

gSetFunction(ctl, fun);
@end group
@end example
@sp 1
See also:  @code{Perform}
@end deffn



@deffn {SetStringValue} SetStringValue::PushButton
@sp 2
@example
@group
r = gSetStringValue(ctl, val);

object  ctl;    /*  push button object  */
char    *val;   /*  display text        */
object  r;      /*  ctl                 */
@end group
@end example
This method sets the display text shown on the button.

The value returned is the control object passed.
@sp 1
See also:  @code{StringValue}
@end deffn



@deffn {StringValue} StringValue::PushButton
@sp 2
@example
@group
val = gStringValue(ctl);

object  ctl;    /*  push button object   */
char    *val;   /*  current display text  */
@end group
@end example
This method returns the button's current display text.
@sp 1
See also:  @code{SetStringValue}
@end deffn















@subsection Check Boxes
The @code{CheckBox} class represents a control which allows the user
to select or de-select a particular option.  Unlike radio buttons,
check boxes are not mutually exclusive.

The @code{CheckBox} class is a subclass of @code{Control} and as such
inherits all the functionality associated with that class.  This section
only documents functionality particular to this class.

Standard control arguments are documented in the section entitled
``Standard Control Method Arguments''.








@deffn {Attach} Attach::CheckBox
@sp 2
@example
@group
r = gAttach(ctl, val);

object  ctl;   /*  control object  */
object  val;   /*  ctl value       */
object  r;     /*  control object  */
@end group
@end example
This method is used to associate an independent Dynace object with the
state associated with a control object.  @code{val} should be an
instance of the @code{ShortInteger} class.  This object will be
automatically updated to reflect the state associated with the control.

The object will represent 0 if the box is not checked, 1 if it is
checked, or 2 if the box is grayed.

This object (@code{val}) will never be disposed by WDS, even after
the control or associated dialog are disposed.  Therefore, this
is one way of gaining access to a control's value after the life
of the control.  It is the programmer's responsibility to dispose of
the object when it is no longer needed.
@example
@group
@exdent Example:

object  ctl, val;

val = gNewWithInt(ShortInteger, 0);
gAttach(ctl, val);
@end group
@end example
@sp 1
See also:  @code{Value, ShortValue, SetValue, SetShortValue}
@end deffn












@deffn {NewCtl} NewCtl::CheckBox
@sp 2
@example
@group
ctl = mNewCtl(CheckBox, id);

unsigned  id;   /*  control id      */
object   ctl;   /*  control object  */
@end group
@end example
@code{mNewCtl} is a macro (defined in @code{dynwin.h}) that expands to
a call to @code{vNew}.

This method is used to create a new control object to be identified as
@code{id} (see section ``Standard Control Method Arguments'').  This
method is mainly used internally.  A programmer would more often
use @code{AddControl::Dialog} to create controls and associate them
to a dialog.
@example
@group
@exdent Example:

object  ctl;

ctl = mNewCtl(CheckBox, MY_CHECKBOX);
@end group
@end example
@sp 1
See also:  @code{AddControl::Dialog, GetControl::Dialog}
@end deffn







@deffn {SetShortValue} SetShortValue::CheckBox
@sp 2
@example
@group
r = gSetShortValue(ctl, val);

object  ctl;    /*  control object  */
int     val;    /*  ctl value       */
object  r;      /*  control object  */
@end group
@end example
This method is used to set the value associated with a control.  It is
often used to set the initial value prior to performing a dialog.
@code{val} should be the value which represents the desired state associated
with the check box.

@code{val} should be set to 0 to uncheck the box, 1 to check it, or 2 to
gray the box (only work on 3-state boxes).

Any previously associated value object (associated via @code{SetValue})
will be disposed when the new value is set.
@example
@group
@exdent Example:

object  ctl;

gSetShortValue(ctl, 1);
@end group
@end example
@sp 1
See also:  @code{SetValue, ShortValue, CtlShortValue::Dialog}
@end deffn
















@deffn {SetValue} SetValue::CheckBox
@sp 2
@example
@group
r = gSetValue(ctl, val);

object  ctl;    /*  control object  */
object  val;    /*  ctl value       */
object  r;      /*  control object  */
@end group
@end example
This method is used to set the value associated with a control.  It is
often used to set the initial value prior to performing a dialog.
@code{val} should be a Dynace object which is an instance of the
@code{ShortInteger} class and initialized to the value desired for the
control.

@code{val} should represent 0 to uncheck the box, 1 to check it, or 2 to
gray the box (only work on 3-state boxes).

Any previously associated object will be disposed when the new value is set.
Also, @code{val} will automatically be disposed when the control or associated
dialog is disposed.
@example
@group
@exdent Example:

object  ctl;

gSetValue(ctl, gNewWithInt(ShortInteger, 1));
@end group
@end example
@sp 1
See also:  @code{SetShortValue, Value}
@end deffn













@deffn {ShortValue} ShortValue::CheckBox
@sp 2
@example
@group
val = gShortValue(ctl);

object  ctl;   /*  control object  */
short   val;   /*  ctl value       */
@end group
@end example
This method is used to obtain a C short integer which represents the
state associated with the control.  The returned integer will be 0 if
the box is not checked, 1 if it is checked, or 2 if the box is grayed.
@example
@group
@exdent Example:

object  ctl;
short   val;

val = gShortValue(ctl);
@end group
@end example
@sp 1
See also:  @code{Value, Attach, SetShortValue, CtlShortValue::Dialog}
@end deffn













@deffn {Value} Value::CheckBox
@sp 2
@example
@group
val = gValue(ctl);

object  ctl;   /*  control object  */
object  val;   /*  ctl value       */
@end group
@end example
This method is used to obtain a Dynace object which represents the value
associated with the control.  The object returned will be an instance of
the @code{ShortInteger} class.  The returned object will represent 0 if
the box is not checked, 1 if it is checked, or 2 if the box is grayed.
This object will be disposed by WDS when the control object or
associated dialog is disposed.
@example
@group
@exdent Example:

object  ctl, val;

val = gValue(ctl);
@end group
@end example
@sp 1
See also:  @code{ShortValue, Attach, SetValue}
@end deffn















@subsection Radio Buttons
The @code{RadioButton} class represents a control which allows the user
to select or de-select a particular option.  Unlike check boxes, however,
radio buttons are mutually exclusive within a group - similar to the
buttons on old fashioned radios, when one button gets selected all the
other ones in the group get de-selected.

When using radio buttons, be sure to use the @code{Group} function in order
to tell WDS which radio buttons should be grouped together.

The @code{RadioButton} class is a subclass of @code{Control} and as such
inherits all the functionality associated with that class.  This section
only documents functionality particular to this class.

Standard control arguments are documented in the section entitled
``Standard Control Method Arguments''.






@deffn {Attach} Attach::RadioButton
@sp 2
@example
@group
r = gAttach(ctl, val);

object  ctl;   /*  control object  */
object  val;   /*  ctl value       */
object  r;     /*  control object  */
@end group
@end example
This method is used to associate an independent Dynace object with the
state associated with a control object.  @code{val} should be an
instance of the @code{ShortInteger} class.  This object will be
automatically updated to reflect the state associated with the control.

The object will represent 0 if the button is not selected, 1 if it is
selected, or 2 if the button is grayed.

This object (@code{val}) will never be disposed by WDS, even after
the control or associated dialog are disposed.  Therefore, this
is one way of gaining access to a control's value after the life
of the control.  It is the programmer's responsibility to dispose of
the object when it is no longer needed.
@example
@group
@exdent Example:

object  ctl, val;

val = gNewWithInt(ShortInteger, 0);
gAttach(ctl, val);
@end group
@end example
@sp 1
See also:  @code{Value, ShortValue, SetValue, SetShortValue}
@end deffn











@deffn {Group} Group::RadioButton
@sp 2
@example
@group
r = mGroup(ctl, id1, id2);

object  ctl;    /*  control object   */
unsigned id1;   /*  starting ctl id  */
unsigned id2;   /*  ending ctl id    */
object  r;      /*  control object   */
@end group
@end example
This method is used to tell WDS which radio buttons are to be considered
part of the same mutually exclusive group.  @code{id1} and @code{id2}
are the resource editor defined id's associated with the beginning and
end of the group.  WDS uses this information to automatically de-select
all radio buttons which are part of a group when one of the buttons is
selected.
@example
@group
@exdent Example:

object  ctl;

mGroup(ctl, FIRST_RB, LAST_RB);
@end group
@end example
@c @sp 1
@c See also:  @code{}
@end deffn



@deffn {GroupShortValue} GroupShortValue::RadioButton
@sp 2
@example
@group
idx = gGroupShortValue(ctl);

object  ctl;   /*  radio button object  */
int     idx;   /*  0-based index or -1  */
@end group
@end example
This method returns the zero-based index of the currently selected
radio button within the group.  Returns -1 if no button is selected.
@sp 1
See also:  @code{GroupStringValue, SetGroupShortValue}
@end deffn



@deffn {GroupStringValue} GroupStringValue::RadioButton
@sp 2
@example
@group
name = gGroupStringValue(ctl);

object  ctl;    /*  radio button object  */
char    *name;  /*  name of selected button  */
@end group
@end example
This method returns the control name of the currently selected radio
button within the group.
@sp 1
See also:  @code{GroupShortValue, SetGroupStringValue}
@end deffn











@deffn {NewCtl} NewCtl::RadioButton
@sp 2
@example
@group
ctl = mNewCtl(RadioButton, id);

unsigned  id;   /*  control id      */
object   ctl;   /*  control object  */
@end group
@end example
@code{mNewCtl} is a macro (defined in @code{dynwin.h}) that expands to
a call to @code{vNew}.

This method is used to create a new control object to be identified as
@code{id} (see section ``Standard Control Method Arguments'').  This
method is mainly used internally.  A programmer would more often
use @code{AddControl::Dialog} to create controls and associate them
to a dialog.
@example
@group
@exdent Example:

object  ctl;

ctl = mNewCtl(RadioButton, MY_RADIOBUTTON);
@end group
@end example
@sp 1
See also:  @code{AddControl::Dialog, GetControl::Dialog, Group}
@end deffn



@deffn {SetGroupShortValue} SetGroupShortValue::RadioButton
@sp 2
@example
@group
r = gSetGroupShortValue(ctl, index);

object  ctl;    /*  radio button object  */
int     index;  /*  0-based index        */
object  r;      /*  ctl                  */
@end group
@end example
This method selects the radio button at the given zero-based index
within the group.

The value returned is the control object passed.
@sp 1
See also:  @code{GroupShortValue, SetGroupStringValue}
@end deffn



@deffn {SetGroupStringValue} SetGroupStringValue::RadioButton
@sp 2
@example
@group
r = gSetGroupStringValue(ctl, name);

object  ctl;    /*  radio button object  */
char    *name;  /*  control name         */
object  r;      /*  ctl                  */
@end group
@end example
This method selects the radio button within the group that has the
specified control name.

The value returned is the control object passed.
@sp 1
See also:  @code{GroupStringValue, SetGroupShortValue}
@end deffn







@deffn {SetShortValue} SetShortValue::RadioButton
@sp 2
@example
@group
r = gSetShortValue(ctl, val);

object  ctl;    /*  control object  */
int     val;    /*  ctl value       */
object  r;      /*  control object  */
@end group
@end example
This method is used to set the value associated with a control.  It is
often used to set the initial value prior to performing a dialog.
@code{val} should be the value which represents the desired state associated
with the radio button.

@code{val} should be set to 0 to de-select the button, 1 to select it, or 2 to
gray the button (only work on 3-state buttons).

Any previously associated value object (associated via @code{SetValue})
will be disposed when the new value is set.
@example
@group
@exdent Example:

object  ctl;

gSetShortValue(ctl, 1);
@end group
@end example
@sp 1
See also:  @code{SetValue, ShortValue, CtlShortValue::Dialog, Group}
@end deffn
















@deffn {SetValue} SetValue::RadioButton
@sp 2
@example
@group
r = gSetValue(ctl, val);

object  ctl;    /*  control object  */
object  val;    /*  ctl value       */
object  r;      /*  control object  */
@end group
@end example
This method is used to set the value associated with a control.  It is
often used to set the initial value prior to performing a dialog.
@code{val} should be a Dynace object which is an instance of the
@code{ShortInteger} class and initialized to the value desired for the
control.

@code{val} should represent 0 to de-select the button, 1 to select it, or 2 to
gray the button (only work on 3-state buttons).

Any previously associated object will be disposed when the new value is set.
Also, @code{val} will automatically be disposed when the control or associated
dialog is disposed.
@example
@group
@exdent Example:

object  ctl;

gSetValue(ctl, gNewWithInt(ShortInteger, 1));
@end group
@end example
@sp 1
See also:  @code{SetShortValue, Value, Group}
@end deffn













@deffn {ShortValue} ShortValue::RadioButton
@sp 2
@example
@group
val = gShortValue(ctl);

object  ctl;   /*  control object  */
short   val;   /*  ctl value       */
@end group
@end example
This method is used to obtain a C short integer which represents the
state associated with the control.  The returned integer will be 0 if
the button is not selected, 1 if it is selected, or 2 if the button is grayed.
@example
@group
@exdent Example:

object  ctl;
short   val;

val = gShortValue(ctl);
@end group
@end example
@sp 1
See also:  @code{Value, Attach, SetShortValue, CtlShortValue::Dialog}
@end deffn













@deffn {Value} Value::RadioButton
@sp 2
@example
@group
val = gValue(ctl);

object  ctl;   /*  control object  */
object  val;   /*  ctl value       */
@end group
@end example
This method is used to obtain a Dynace object which represents the value
associated with the control.  The object returned will be an instance of
the @code{ShortInteger} class.  The returned object will represent 0 if
the button is not selected, 1 if it is selected, or 2 if the button is grayed.
This object will be disposed by WDS when the control object or
associated dialog is disposed.
@example
@group
@exdent Example:

object  ctl, val;

val = gValue(ctl);
@end group
@end example
@sp 1
See also:  @code{ShortValue, Attach, SetValue}
@end deffn











@subsection List Boxes
The @code{ListBox} class represents a control which allows the user
to select from a list of choices.  These choices are presented in
a vertical, possibly drop down, list of text choices.  This control
supports access to the user's choice via the text option selected or
an ordinal value.

The application specifies, at runtime, the choices located in the list box.

The main difference between list boxes and combo boxes is that with list
boxes the user can only select from the available choices, however, in
addition, combo boxes allow the user to enter a choice independent of
the available options.

The @code{ListBox} class is a subclass of @code{Control} and as such
inherits all the functionality associated with that class.  This section
only documents functionality particular to this class.

Standard control arguments are documented in the section entitled
``Standard Control Method Arguments''.













@deffn {AddOption} AddOption::ListBox
@sp 2
@example
@group
r = gAddOption(ctl, op);

object  ctl;    /*  control object  */
char    *op;    /*  new option      */
object  r;      /*  control object  */
@end group
@end example
This method is used to append a new option to the list of options associated
with a list box.  This method may be used prior to or during the execution
(via @code{gPerform}) of a dialog.

@code{op} may also be a Dynace object, an instance of the @code{String}
class, representing the new choice.  If a Dynace object is used, it must
be typecast to a string and it will be disposed when the control or
associated dialog are disposed.
@example
@group
@exdent Example:

object  ctl;

gAddOption(ctl, "The first choice");
gAddOption(ctl, "The second choice");
@end group
@end example
@sp 1
See also:  @code{Alphabetize::ListBox, AddOptionAt::ListBox}
@end deffn



@deffn {AddOptionAt} AddOptionAt::ListBox
@sp 2
@example
@group
r = gAddOptionAt(ctl, pos, ops);

object  ctl;    /*  list box object   */
int     pos;    /*  insertion position (0-based)  */
char    *ops;   /*  option text       */
object  r;      /*  ctl               */
@end group
@end example
This method inserts an option at position @code{pos} in the list box.

The value returned is the control object passed.
@sp 1
See also:  @code{AddOption, RemoveAll}
@end deffn















@deffn {Alphabetize} Alphabetize::ListBox
@sp 2
@example
@group
r = gAlphabetize(ctl);

object  ctl;    /*  control object  */
object  r;      /*  control object  */
@end group
@end example
This method is used to cause the automatic alphabetization of all
options added via the @code{AddOption} method.  In order to function
properly, @code{Alphabetize} must be called prior to any calls to
@code{AddOption}.  This method simply returns the object passed.
@c @example
@c @group
@c @exdent Example:
@c 
@c @end group
@c @end example
@sp 1
See also:  @code{AddOption::ListBox}
@end deffn















@deffn {Attach} Attach::ListBox
@sp 2
@example
@group
r = gAttach(ctl, val);

object  ctl;   /*  control object  */
object  val;   /*  ctl value       */
object  r;     /*  control object  */
@end group
@end example
This method is used to associate an independent Dynace object with the
state associated with a control object.  @code{val} should be an
instance of the @code{ShortInteger} class or the @code{String} class.
This object will be automatically updated to reflect the state
associated with the control.

If @code{val} is an instance of the @code{ShortInteger} class it will
be used to represent the ordinal value of the selected item.  This value
is a 0 based index from the top of the list to the bottom.  -1 is used
to represent no valid selection.

If @code{val} is an instance of the @code{String} class it will be used
to represent the string represented by the user's choice.  If no choice is
made it will represent "".

This object (@code{val}) will never be disposed by WDS, even after
the control or associated dialog are disposed.  Therefore, this
is one way of gaining access to a control's value after the life
of the control.  It is the programmer's responsibility to dispose of
the object when it is no longer needed.
@example
@group
@exdent Example:

object  ctl, val;

val = gNewWithInt(ShortInteger, 0);
gAttach(ctl, val);
@end group
@end example
@sp 1
See also:  @code{Value, ShortValue, StringValue}
@end deffn



@deffn {DeselectLine} DeselectLine::ListBox
@sp 2
@example
@group
r = gDeselectLine(ctl, line);

object  ctl;    /*  list box object  */
int     line;   /*  0-based line index  */
object  r;      /*  ctl              */
@end group
@end example
This method deselects the specified line in a multi-selection list box.
@code{line} is a zero-based index.

The value returned is the control object passed.
@sp 1
See also:  @code{SelectLine, GetSelections}
@end deffn
















@deffn {GetSelections} GetSelections::ListBox
@sp 2
@example
@group
ary = gGetSelections(ctl);

object  ctl;   /*  control object       */
object  ary;   /*  array of selections  */
@end group
@end example
This method is used to obtain an array representing the items selected
by the user.  It is most useful with list boxes which have multi-selection
enabled.  The array returned is an instance of the @code{IntegerArray} class
and @emph{must} be explicitly disposed when no longer needed.  Each element
of the array represents a zero based index of one of the user's selections,
and the size of the array represents the number of items selected.

This method may only be used while a dialog is active and will return
@code{NULL} otherwise.
@c @example
@c @group
@c @exdent Example:
@c 
@c @end group
@c @end example
@sp 1
See also:  @code{ValueAt, NumbSelected, Value, ListIndex}
@end deffn
















@deffn {ListIndex} ListIndex::ListBox
@sp 2
@example
@group
val = gListIndex(ctl);

object  ctl;   /*  control object  */
object  val;   /*  ctl value       */
@end group
@end example
This method is used to obtain a Dynace object which represents the value
associated with the control.  The object returned will be an instance of
the @code{ShortInteger} class.  The returned object will represent the
ordinal value of the choice made by the user.  This ordinal value is a
zero based index from top to bottom.  -1 indicates that no choice was
made.

This object will be disposed by WDS when the control object or
associated dialog is disposed.
@example
@group
@exdent Example:

object  ctl, val;

val = gListIndex(ctl);
@end group
@end example
@sp 1
See also:  @code{ShortValue, Value, StringValue, GetSelections, SetValue}
@end deffn




@deffn {NewCtl} NewCtl::ListBox
@sp 2
@example
@group
ctl = mNewCtl(ListBox, id);

unsigned  id;   /*  control id      */
object   ctl;   /*  control object  */
@end group
@end example
@code{mNewCtl} is a macro (defined in @code{dynwin.h}) that expands to
a call to @code{vNew}.

This method is used to create a new control object to be identified as
@code{id} (see section ``Standard Control Method Arguments'').  This
method is mainly used internally.  A programmer would more often
use @code{AddControl::Dialog} to create controls and associate them
to a dialog.
@example
@group
@exdent Example:

object  ctl;

ctl = mNewCtl(ListBox, MY_LISTBOX);
@end group
@end example
@sp 1
See also:  @code{AddControl::Dialog, GetControl::Dialog, AddOption}
@end deffn











@deffn {NumbSelected} NumbSelected::ListBox
@sp 2
@example
@group
n = gNumbSelected(ctl);

object  ctl;    /*  control object   */
int     n;      /*  number selected  */
@end group
@end example
This method is used to obtain the number of list box items the user
has selected.  It is mainly used with multi-selection list boxes.
This method should only be used while the dialog is active and
will return -1 otherwise.
@c @example
@c @group
@c @exdent Example:
@c 
@c @end group
@c @end example
@sp 1
See also:  @code{GetSelections}
@end deffn














@deffn {RemoveAll} RemoveAll::ListBox
@sp 2
@example
@group
r = gRemoveAll(ctl);

object  ctl;    /*  control object  */
object  r;      /*  control object  */
@end group
@end example
This method is used to remove all items from a list box.  It may be used
prior to or during the execution of a dialog.  The control object passed
is returned.
@c @example
@c @group
@c @exdent Example:
@c 
@c @end group
@c @end example
@sp 1
See also:  @code{RemoveStr, RemoveInt}
@end deffn
















@deffn {RemoveInt} RemoveInt::ListBox
@sp 2
@example
@group
r = gRemoveInt(ctl, idx);

object  ctl;    /*  control object  */
int     idx;    /*  index           */
object  r;      /*  control object  */
@end group
@end example
This method is used to remove an item from a list box while the dialog
is active.  @code{idx} is a zero based index of the item to be removed.
If the operation succeeded @code{ctl} is returned, otherwise @code{NULL}
is returned.

This method may only be used while a dialog is active.  It should not be
used prior to performing (via @code{gPerform}) a dialog or after the user
has accepted or canceled the dialog.
@example
@group
@exdent Example:

object  ctl;

gRemoveInt(ctl, 1);
@end group
@end example
@sp 1
See also:  @code{RemoveStr, RemoveAll}
@end deffn














@deffn {RemoveStr} RemoveStr::ListBox
@sp 2
@example
@group
r = gRemoveStr(ctl, itm);

object  ctl;    /*  control object  */
char   *itm;    /*  item            */
object  r;      /*  control object  */
@end group
@end example
This method is used to remove an item from a list box while the dialog
is active.  @code{itm} is the string to be removed.
If the operation succeeded @code{ctl} is returned, otherwise @code{NULL}
is returned.

This method may only be used while a dialog is active.  It should not be
used prior to performing (@code{gPerform}) a dialog or after the user
has accepted or canceled the dialog.
@example
@group
@exdent Example:

object  ctl;

gRemoveStr(ctl, "Some Option");
@end group
@end example
@sp 1
See also:  @code{RemoveInt, FindMode}
@end deffn















@deffn {Required} Required::ListBox
@sp 2
@example
@group
r = gRequired(ctl, mode);

object  ctl;    /*  control object  */
int     mode;   /*  required mode   */
object  r;      /*  control object  */
@end group
@end example
This method is used to determine whether or not the user is required to
make a valid selection prior to WDS allowing the acceptance of the dialog.
If it is required and the user does not make a valid selection, WDS
will issue an error message and return them to the control.

@code{mode} is 1 to make it required and 0 otherwise.
@example
@group
@exdent Example:

object  ctl;

gRequired(ctl, 1);
@end group
@end example
@sp 1
See also:  @code{CheckFunction::Control, CheckValue::Control}
@end deffn



@deffn {SelectLine} SelectLine::ListBox
@sp 2
@example
@group
r = gSelectLine(ctl, line);

object  ctl;    /*  list box object  */
int     line;   /*  0-based line index  */
object  r;      /*  ctl              */
@end group
@end example
This method selects the specified line in a multi-selection list box.
@code{line} is a zero-based index.

The value returned is the control object passed.
@sp 1
See also:  @code{DeselectLine, GetSelections}
@end deffn








@deffn {SetFunction} SetFunction::ListBox
@sp 2
@example
@group
r = gSetFunction(ctl, fun);

object  ctl;      /*  control object           */
int     (*fun)(); /*  function pointer         */
ofun    r;        /*  previous function        */
@end group
@end example
This method is inherited from the @code{Control} class.

This method is used to associate a C function to a list box such that if
the user double clicks an item in the list box the C function will
immediately get evoked.  The value returned by this method is the
function, if any, which was previously associated with the control.

The C function takes the following form:
@example
@group
int     fun(object ctl, object dlg)
@{
        ....
@}
@end group
@end example
Where @code{ctl} is the control object, and @code{dlg} is the object
representing the dialog the control is associated with.  The value
returned by @code{fun} is ignored.
@example
@group
@exdent Example:

object  ctl;

gSetFunction(ctl, fun);
@end group
@end example
@sp 1
See also:  @code{Perform, SetChgFunction}
@end deffn


















@deffn {SetShortValue} SetShortValue::ListBox
@sp 2
@example
@group
r = gSetShortValue(ctl, val);

object  ctl;    /*  control object  */
int     val;    /*  ctl value       */
object  r;      /*  control object  */
@end group
@end example
This method is used to set the default selection associated with the
control.  It is often used to set the initial value prior to performing
a dialog.

@code{val} will be used as the ordinal value of the selected item.  This
value is a 0 based index from the top of the list to the bottom.

Any previously associated object (via @code{Value}) will be disposed
when the new value is set.
@example
@group
@exdent Example:

object  ctl;

gSetShortValue(ctl, 1);
@end group
@end example
@sp 1
See also:  @code{SetStringValue, SetValue}
@end deffn









@deffn {SetStringValue} SetStringValue::ListBox
@sp 2
@example
@group
r = gSetStringValue(ctl, val);

object  ctl;    /*  control object  */
char    *val;   /*  ctl value       */
object  r;      /*  control object  */
@end group
@end example
This method is used to set the default selection associated with the
control.  It is often used to set the initial value prior to performing
a dialog.  

@code{val} will be used to select the option which matches the string.

Any previously associated object (via @code{val}) will be disposed when
the new value is set.  
@example
@group
@exdent Example:

object  ctl;

gSetStringValue(ctl, "The choice");
@end group
@end example
@sp 1
See also:  @code{SetShortValue, SetValue}
@end deffn










@deffn {SetValue} SetValue::ListBox
@sp 2
@example
@group
r = gSetValue(ctl, val);

object  ctl;    /*  control object  */
object  val;    /*  ctl value       */
object  r;      /*  control object  */
@end group
@end example
This method is used to set the default selection associated with the
control.  It is often used to set the initial value prior to performing
a dialog.  @code{val} should be a Dynace object which is an instance of
either the @code{String} class or the @code{ShortInteger} class and
initialized to the value desired for the control.

If @code{val} is an instance of the @code{ShortInteger} class it will be
used as the ordinal value of the selected item.  This value is a 0 based
index from the top of the list to the bottom.

If @code{val} is an instance of the @code{String} class it will be used
to select the option which matches the string represented by @code{val}.

Any previously associated object will be disposed when the new value is set.
Also, @code{val} will automatically be disposed when the control or associated
dialog is disposed.
@example
@group
@exdent Example:

object  ctl;

gSetValue(ctl, gNewWithInt(ShortInteger, 1));
        or
gSetValue(ctl, gNewWithStr(String, "The choice"));
@end group
@end example
@sp 1
See also:  @code{SetShortValue, SetStringValue}
@end deffn













@deffn {ShortValue} ShortValue::ListBox
@sp 2
@example
@group
val = gShortValue(ctl);

object  ctl;   /*  control object  */
short   val;   /*  ctl value       */
@end group
@end example
This method is used to obtain a C short which represents the value
associated with the control.  The returned value will represent the
ordinal value of the choice made by the user.  This ordinal value is a
zero based index from top to bottom.  -1 indicates that no choice was
made.
@example
@group
@exdent Example:

object  ctl;
short   val;

val = gShortValue(ctl);
@end group
@end example
@sp 1
See also:  @code{ListIndex, Value, GetSelections, StringValue, SetValue}
@end deffn



@deffn {Size} Size::ListBox
@sp 2
@example
@group
n = gSize(ctl);

object  ctl;   /*  list box object   */
int     n;     /*  number of items   */
@end group
@end example
This method returns the number of items in the list box.
@end deffn










@deffn {StringValue} StringValue::ListBox
@sp 2
@example
@group
val = gStringValue(ctl);

object  ctl;   /*  control object  */
char    *val;  /*  ctl value       */
@end group
@end example
This method is used to obtain a character string pointer which
represents the value associated with the control.  The returned pointer
will represent the text associated with the user selected choices.
It will point to "" if the user had not made a valid selection.

The returned pointer will not be valid once the control object or
associated dialog is disposed.
@example
@group
@exdent Example:

object  ctl;
char    *val;

val = gStringValue(ctl);
@end group
@end example
@sp 1
See also:  @code{ShortValue, Value, ListIndex, GetSelections, SetValue}
@end deffn










@deffn {Value} Value::ListBox
@sp 2
@example
@group
val = gValue(ctl);

object  ctl;   /*  control object  */
object  val;   /*  ctl value       */
@end group
@end example
This method is used to obtain a Dynace object which represents the value
associated with the control.  The object returned will be an instance of
the @code{String} class.  The returned object will represent the text
associated with the user selected choices.  If no selection was made
the object will represent "".

This object will be disposed by WDS when the control object or
associated dialog is disposed.
@example
@group
@exdent Example:

object  ctl, val;

val = gValue(ctl);
@end group
@end example
@sp 1
See also:  @code{ShortValue, StringValue, ValueAt, ListIndex, GetSelections, SetValue}
@end deffn









@deffn {ValueAt} ValueAt::ListBox
@sp 2
@example
@group
val = gValueAt(ctl, idx);

object  ctl;   /*  control object  */
int     idx;   /*  index           */
object  val;   /*  value at idx    */
@end group
@end example
This method is used to obtain a Dynace object which represents the text
associated with the line indexed by the zero based index @code{idx}.
The object returned will be an instance of the @code{String} class.  If
the index is out of range this method will return @code{NULL}.

This method may only be called when the dialog is active.  The object
returned @emph{must} be explicitly disposed when no longer needed.
@c @example
@c @group
@c @exdent Example:
@c
@c @end group
@c @end example
@sp 1
See also:  @code{ShortValue, Value, ListIndex, GetSelections, SetValue}
@end deffn








@subsection Combo Boxes
The @code{ComboBox} class represents a control which allows the user to
select from a list of choices or enter an alternate text string.  These
choices are presented in a vertical, possibly drop down, list of text
choices.  This control supports access to the user's choice via the text
option selected or an ordinal value.

The application specifies, at runtime, the choices located in the list box.

The main difference between list boxes and combo boxes is that with list
boxes the user can only select from the available choices, however, in
addition, combo boxes allow the user to enter a choice independent of
the available options.

The @code{ComboBox} class is a subclass of @code{Control} and as such
inherits all the functionality associated with that class.  This section
only documents functionality particular to this class.

Standard control arguments are documented in the section entitled
``Standard Control Method Arguments''.










@deffn {AddEdtHandlerAfter} AddEdtHandlerAfter::ComboBox
@sp 2
@example
@group
r = gAddEdtHandlerAfter(ctl, msg, func);

object   ctl;      /*  a control object  */
unsigned msg;      /*  message           */
long    (*func)(); /*  function pointer  */
object  r;         /*  the control obj   */
@end group
@end example
This method is used to associate function @code{func} with the text
entry window portion of the combo box, message @code{msg} for
control @code{ctl}.  Whenever the text entry window of control
@code{ctl} receives message @code{msg}, @code{func} will be called.

@code{ctl} is the control object who's text window messages you wish to
process.  @code{msg} is the particular message you wish to trap.  These
messages are fully documented in the Windows documentation in the
Messages section.  They normally begin with @code{WM_}.

@code{func} is the function which gets called whenever the specified
message gets received and takes the following form:
@example
@group
long    func(object     ctl,
             HWND       hwnd, 
             UINT       mMsg, 
             WPARAM     wParam, 
             LPARAM     lParam)
@{
        .
        .
        .
        return 0L;  /* or whatever is appropriate  */
@}
@end group
@end example
Where @code{ctl} is the control being sent the message.  The remaining
arguments and return value is fully documented in the Windows documentation
under the @code{WindowProc} function and the Windows Messages documentation.

WDS keeps a list of functions associated with each message associated
with each control.  When a particular message is received the appropriate
list of handler functions gets executed sequentially.
@code{AddEdtHandlerAfter} appends the new function to the end of this list,
and @code{AddEdtHandlerBefore} adds the new function to the beginning of
the list.

WDS may also, and optionally, execute the Windows default procedure
associated with a given message either before or after the user added
list of functions.  This behavior may be controlled via
@code{DefaultEdtProcessingMode}.

Windows will only see the return value of the last message handler executed
including, if applicable, the default.
@example
@group
@exdent Example:

int     hSize, vSize;

static  long    process_wm_size(object  ctl, 
                                HWND    hwnd, 
                                UINT    mMsg, 
                                WPARAM  wParam, 
                                LPARAM  lParam)
@{
        hSize = LOWORD(lParam);
        vSize = HIWORD(lParam);
        return 0L;
@}

        .
        .
        gAddEdtHandlerAfter(ctl, (unsigned) WM_SIZE,
                                       process_wm_size);
        .
        .
@end group
@end example
@sp 1
See also:  @code{DefaultProcessingMode, AddEdtHandlerBefore,}
@iftex
@hfil @break @hglue .63in 
@end iftex
@code{AddHandlerAfter::Control}
@end deffn






@deffn {AddEdtHandlerBefore} AddEdtHandlerBefore::ComboBox
@sp 2
@example
@group
r = gAddEdtHandlerBefore(ctl, msg, func);

object   ctl;      /*  a control object  */
unsigned msg;      /*  message           */
long    (*func)(); /*  function pointer  */
object  r;         /*  the control obj   */
@end group
@end example
This function is fully documented under @code{AddEdtHandlerAfter}.
@c @example
@c @group
@c @exdent Example:
@c @end group
@c @end example
@sp 1
See also:  @code{AddEdtHandlerAfter, AddHandlerAfter::Control}
@end deffn













@deffn {AddOption} AddOption::ComboBox
@sp 2
@example
@group
r = gAddOption(ctl, op);

object  ctl;    /*  control object  */
char    *op;    /*  new option      */
object  r;      /*  control object  */
@end group
@end example
This method is used to append a new option to the list of options associated
with a combo box.  This method may be used prior to or during the execution
(via @code{gPerform}) of a dialog.

@code{op} may also be a Dynace object, an instance of the @code{String}
class, representing the new choice.  If a Dynace object is used, it must
be typecast to a string and it will be disposed when the control or
associated dialog are disposed.
@example
@group
@exdent Example:

object  ctl;

gAddOption(ctl, "The first choice");
gAddOption(ctl, "The second choice");
@end group
@end example
@sp 1
See also:  @code{Alphabetize::ComboBox, AddOptionAt::ComboBox}
@end deffn



@deffn {AddOptionAt} AddOptionAt::ComboBox
@sp 2
@example
@group
r = gAddOptionAt(ctl, pos, ops);

object  ctl;    /*  combo box object  */
int     pos;    /*  insertion position (0-based)  */
char    *ops;   /*  option text       */
object  r;      /*  ctl               */
@end group
@end example
This method inserts an option at position @code{pos} in the combo box.

The value returned is the control object passed.
@sp 1
See also:  @code{AddOption, RemoveAll}
@end deffn












@deffn {Alphabetize} Alphabetize::ComboBox
@sp 2
@example
@group
r = gAlphabetize(ctl);

object  ctl;    /*  control object  */
object  r;      /*  control object  */
@end group
@end example
This method is used to cause the automatic alphabetization of all
options added via the @code{AddOption} method.  In order to function
properly, @code{Alphabetize} must be called prior to any calls to
@code{AddOption}.  This method simply returns the object passed.
@c @example
@c @group
@c @exdent Example:
@c 
@c @end group
@c @end example
@sp 1
See also:  @code{AddOption::ComboBox}
@end deffn












@deffn {Attach} Attach::ComboBox
@sp 2
@example
@group
r = gAttach(ctl, val);

object  ctl;   /*  control object  */
object  val;   /*  ctl value       */
object  r;     /*  control object  */
@end group
@end example
This method is used to associate an independent Dynace object with the
state associated with a control object.  @code{val} should be an
instance of the @code{ShortInteger} class or the @code{String} class.
This object will be automatically updated to reflect the state
associated with the control.

If @code{val} is an instance of the @code{ShortInteger} class it will
be used to represent the ordinal value of the selected item.  This value
is a 0 based index from the top of the list to the bottom.  -1 is used
to represent no valid selection.

If @code{val} is an instance of the @code{String} class it will be used
to represent the string represented by the user's choice.  If no choice is
made it will represent "".

This object (@code{val}) will never be disposed by WDS, even after
the control or associated dialog are disposed.  Therefore, this
is one way of gaining access to a control's value after the life
of the control.  It is the programmer's responsibility to dispose of
the object when it is no longer needed.
@example
@group
@exdent Example:

object  ctl, val;

val = gNewWithInt(ShortInteger, 0);
gAttach(ctl, val);
@end group
@end example
@sp 1
See also:  @code{Value, ShortValue, StringValue}
@end deffn












@deffn {DefaultEdtProcessingMode} DefaultEdtProcessingMode::ComboBox
@sp 2
@example
@group
r = gDefaultEdtProcessingMode(ctl, msg, mode);

object   ctl;   /*  a control object         */
unsigned msg;   /*  message                  */
int      mode;  /*  default processing mode  */
object   r;     /*  the control obj          */
@end group
@end example
This method is used to determine when or if the Windows default message
procedure is processed for a given message (@code{msg}) associated with
the text entry window portion of control (@code{ctl}).

WDS allows a programmer to specify an arbitrary number of functions to
be executed whenever the text entry portion of a combo box control
receives a specific message (via @code{AddEdtHandlerAfter} and
@code{AddEdtHandlerBefore}).  Windows has default procedures associated
with many control messages.  At times it is necessary to replace or
augment this default functionality.  @code{DefaultEdtProcessingMode}
gives the programmer control over when and if this default Windows
functionality.  @code{mode} is used to specify the desired mode.  The
following table indicates the valid modes:

@table @code
@item 0
Do not execute the Windows default processing
@item 1
Execute default processing @emph{after} programmer defined handlers
@item 2
Execute default processing @emph{before} programmer defined handlers
@end table

Note that the default mode is always @code{1}, and must be explicitly
changed, if desired, for each message associated with each control.

@code{msg} is the particular message you wish to affect.  These messages
are fully documented in the Windows documentation in the Messages
section.  They normally begin with @code{WM_}.
@example
@group
@exdent Example:

gDefaultEdtProcessingMode(ctl, (unsigned) WM_SIZE, 0);
@end group
@end example
@sp 1
See also:  @code{AddEdtHandlerAfter}
@end deffn



@deffn {FindIndex} FindIndex::ComboBox
@sp 2
@example
@group
idx = gFindIndex(ctl, choice);

object  ctl;      /*  combo box object  */
char    *choice;  /*  text to find      */
int     idx;      /*  0-based index or -1  */
@end group
@end example
This method returns the zero-based index of @code{choice} in the
combo box's option list, or -1 if the choice is not found.
@end deffn

















@deffn {ListIndex} ListIndex::ComboBox
@sp 2
@example
@group
val = gListIndex(ctl);

object  ctl;   /*  control object  */
object  val;   /*  ctl value       */
@end group
@end example
This method is used to obtain a Dynace object which represents the value
associated with the control.  The object returned will be an instance of
the @code{ShortInteger} class.  The returned object will represent the
ordinal value of the choice made by the user.  This ordinal value is a
zero based index from top to bottom.  -1 indicates that no choice was
made.

This object will be disposed by WDS when the control object or
associated dialog is disposed.
@example
@group
@exdent Example:

object  ctl, val;

val = gListIndex(ctl);
@end group
@end example
@sp 1
See also:  @code{ShortValue, Value, StringValue, Attach, SetValue}
@end deffn




@deffn {NewCtl} NewCtl::ComboBox
@sp 2
@example
@group
ctl = mNewCtl(ComboBox, id);

unsigned  id;   /*  control id      */
object   ctl;   /*  control object  */
@end group
@end example
@code{mNewCtl} is a macro (defined in @code{dynwin.h}) that expands to
a call to @code{vNew}.

This method is used to create a new control object to be identified as
@code{id} (see section ``Standard Control Method Arguments'').  This
method is mainly used internally.  A programmer would more often
use @code{AddControl::Dialog} to create controls and associate them
to a dialog.
@example
@group
@exdent Example:

object  ctl;

ctl = mNewCtl(ComboBox, MY_COMBOBOX);
@end group
@end example
@sp 1
See also:  @code{AddControl::Dialog, GetControl::Dialog, AddOption}
@end deffn



@deffn {NumbSelected} NumbSelected::ComboBox
@sp 2
@example
@group
n = gNumbSelected(ctl);

object  ctl;   /*  combo box object  */
int     n;     /*  1 if selected, 0 otherwise  */
@end group
@end example
This method returns 1 if an item is currently selected in the combo
box, or 0 if no item is selected.
@end deffn




















@deffn {RemoveAll} RemoveAll::ComboBox
@sp 2
@example
@group
r = gRemoveAll(ctl);

object  ctl;    /*  control object  */
object  r;      /*  control object  */
@end group
@end example
This method is used to remove all items from a combo box.  It may be used
prior to or during the execution of a dialog.  The control object passed
is returned.
@c @example
@c @group
@c @exdent Example:
@c 
@c @end group
@c @end example
@sp 1
See also:  @code{RemoveStr, RemoveInt}
@end deffn

















@deffn {RemoveInt} RemoveInt::ComboBox
@sp 2
@example
@group
r = gRemoveInt(ctl, idx);

object  ctl;    /*  control object  */
int     idx;    /*  index           */
object  r;      /*  control object  */
@end group
@end example
This method is used to remove an item from a combo box while the dialog
is active.  @code{idx} is a zero based index of the item to be removed.
If the operation succeeded @code{ctl} is returned, otherwise @code{NULL}
is returned.

This method may only be used while a dialog is active.  It should not be
used prior to performing (@code{gPerform}) a dialog or after the user
has accepted or canceled the dialog.
@example
@group
@exdent Example:

object  ctl;

gRemoveInt(ctl, 1);
@end group
@end example
@sp 1
See also:  @code{RemoveStr, RemoveAll}
@end deffn














@deffn {RemoveStr} RemoveStr::ComboBox
@sp 2
@example
@group
r = gRemoveStr(ctl, itm);

object  ctl;    /*  control object  */
char   *itm;    /*  item            */
object  r;      /*  control object  */
@end group
@end example
This method is used to remove an item from a combo box while the dialog
is active.  @code{itm} is the string to be removed.
If the operation succeeded @code{ctl} is returned, otherwise @code{NULL}
is returned.

This method may only be used while a dialog is active.  It should not be
used prior to performing (@code{gPerform}) a dialog or after the user
has accepted or canceled the dialog.
@example
@group
@exdent Example:

object  ctl;

gRemoveStr(ctl, "Some Option");
@end group
@end example
@sp 1
See also:  @code{RemoveInt::ComboBox, FindMode::ComboBox}
@end deffn

















@deffn {Required} Required::ComboBox
@sp 2
@example
@group
r = gRequired(ctl, mode);

object  ctl;    /*  control object  */
int     mode;   /*  required mode   */
object  r;      /*  control object  */
@end group
@end example
This method is used to determine whether or not the user is required to
make a valid selection prior to WDS allowing the acceptance of the dialog.
If it is required and the user does not make a valid selection, WDS
will issue an error message and return them to the control.

@code{mode} is 1 to make it required and 0 otherwise.
@example
@group
@exdent Example:

object  ctl;

gRequired(ctl, 1);
@end group
@end example
@sp 1
See also:  @code{CheckFunction::Control, CheckValue::Control}
@end deffn



@deffn {SetChgFunction} SetChgFunction::ComboBox
@sp 2
@example
@group
ofun = gSetChgFunction(ctl, fun);

object  ctl;         /*  combo box object     */
int     (*fun)();    /*  change callback      */
int     (*ofun)();   /*  previous callback    */
@end group
@end example
This method sets the function to be called when the user changes the
selection in the combo box.

The value returned is the previous change function.
@sp 1
See also:  @code{SetFunction}
@end deffn



















@deffn {SetFunction} SetFunction::ComboBox
@sp 2
@example
@group
r = gSetFunction(ctl, fun);

object  ctl;      /*  control object           */
int     (*fun)(); /*  function pointer         */
ofun    r;        /*  previous function        */
@end group
@end example
This method is inherited from the @code{Control} class.

This method is used to associate a C function to a combo box such that if
the user double clicks an item in the combo box the C function will
immediately get evoked.  The value returned by this method is the
function, if any, which was previously associated with the control.

The C function takes the following form:
@example
@group
int     fun(object ctl, object dlg)
@{
        ....
@}
@end group
@end example
Where @code{ctl} is the control object, and @code{dlg} is the object
representing the dialog the control is associated with.  The value
returned by @code{fun} is ignored.
@example
@group
@exdent Example:

object  ctl;

gSetFunction(ctl, fun);
@end group
@end example
@sp 1
See also:  @code{Perform, SetChgFunction}
@end deffn






















@deffn {SetShortValue} SetShortValue::ComboBox
@sp 2
@example
@group
r = gSetShortValue(ctl, val);

object  ctl;    /*  control object  */
int     val;    /*  ctl value       */
object  r;      /*  control object  */
@end group
@end example
This method is used to set the default selection associated with the
control.  It is often used to set the initial value prior to performing
a dialog.

@code{val} will be used as the ordinal value of the selected item.  This
value is a 0 based index from the top of the list to the bottom.

Any previously associated object (via @code{Value}) will be disposed
when the new value is set.
@example
@group
@exdent Example:

object  ctl;

gSetShortValue(ctl, 1);
@end group
@end example
@sp 1
See also:  @code{SetStringValue, SetValue}
@end deffn









@deffn {SetStringValue} SetStringValue::ComboBox
@sp 2
@example
@group
r = gSetStringValue(ctl, val);

object  ctl;    /*  control object  */
char    *val;   /*  ctl value       */
object  r;      /*  control object  */
@end group
@end example
This method is used to set the default selection associated with the
control.  It is often used to set the initial value prior to performing
a dialog.  

@code{val} will be used to select the option which matches the string.

Any previously associated object (via @code{val}) will be disposed when
the new value is set.  
@example
@group
@exdent Example:

object  ctl;

gSetStringValue(ctl, "The choice");
@end group
@end example
@sp 1
See also:  @code{SetShortValue, SetValue}
@end deffn










@deffn {SetValue} SetValue::ComboBox
@sp 2
@example
@group
r = gSetValue(ctl, val);

object  ctl;    /*  control object  */
object  val;    /*  ctl value       */
object  r;      /*  control object  */
@end group
@end example
This method is used to set the default selection associated with the
control.  It is often used to set the initial value prior to performing
a dialog.  @code{val} should be a Dynace object which is an instance of
either the @code{String} class or the @code{ShortInteger} class and
initialized to the value desired for the control.

If @code{val} is an instance of the @code{ShortInteger} class it will be
used as the ordinal value of the selected item.  This value is a 0 based
index from the top of the list to the bottom.

If @code{val} is an instance of the @code{String} class it will be used
to select the option which matches the string represented by @code{val}.

Any previously associated object will be disposed when the new value is set.
Also, @code{val} will automatically be disposed when the control or associated
dialog is disposed.
@example
@group
@exdent Example:

object  ctl;

gSetValue(ctl, gNewWithInt(ShortInteger, 1));
        or
gSetValue(ctl, gNewWithStr(String, "The choice"));
@end group
@end example
@sp 1
See also:  @code{SetShortValue, SetStringValue}
@end deffn













@deffn {ShortValue} ShortValue::ComboBox
@sp 2
@example
@group
val = gShortValue(ctl);

object  ctl;   /*  control object  */
short   val;   /*  ctl value       */
@end group
@end example
This method is used to obtain a C short which represents the value
associated with the control.  The returned value will represent the
ordinal value of the choice made by the user.  This ordinal value is a
zero based index from top to bottom.  -1 indicates that no choice was
made.
@example
@group
@exdent Example:

object  ctl;
short   val;

val = gShortValue(ctl);
@end group
@end example
@sp 1
See also:  @code{ListIndex, Value, StringValue, Attach, SetValue}
@end deffn



@deffn {Size} Size::ComboBox
@sp 2
@example
@group
n = gSize(ctl);

object  ctl;   /*  combo box object  */
int     n;     /*  number of items   */
@end group
@end example
This method returns the number of items in the combo box.
@end deffn










@deffn {StringValue} StringValue::ComboBox
@sp 2
@example
@group
val = gStringValue(ctl);

object  ctl;   /*  control object  */
char    *val;  /*  ctl value       */
@end group
@end example
This method is used to obtain a character string pointer which
represents the value associated with the control.  The returned pointer
will represent the text associated with the user selected choices.
It will point to "" if the user had not made a valid selection.

The returned pointer will not be valid once the control object or
associated dialog is disposed.
@example
@group
@exdent Example:

object  ctl;
char    *val;

val = gStringValue(ctl);
@end group
@end example
@sp 1
See also:  @code{ShortValue, Value, ListIndex, Attach, SetValue}
@end deffn










@deffn {Value} Value::ComboBox
@sp 2
@example
@group
val = gValue(ctl);

object  ctl;   /*  control object  */
object  val;   /*  ctl value       */
@end group
@end example
This method is used to obtain a Dynace object which represents the value
associated with the control.  The object returned will be an instance of
the @code{String} class.  The returned object will represent the text
associated with the user selected choices.  If no selection was made
the object will represent "".

This object will be disposed by WDS when the control object or
associated dialog is disposed.
@example
@group
@exdent Example:

object  ctl, val;

val = gValue(ctl);
@end group
@end example
@sp 1
See also:  @code{ShortValue, StringValue, ValueAt, ListIndex, Attach, SetValue}
@end deffn






@deffn {ValueAt} ValueAt::ComboBox
@sp 2
@example
@group
val = gValueAt(ctl, idx);

object  ctl;   /*  control object  */
int     idx;   /*  index           */
object  val;   /*  value at idx    */
@end group
@end example
This method is used to obtain a Dynace object which represents the text
associated with the line indexed by the zero based index @code{idx}.
The object returned will be an instance of the @code{String} class.  If
the index is out of range this method will return @code{NULL}.

This method may only be called when the dialog is active.  The object
returned @emph{must} be explicitly disposed when no longer needed.
@c @example
@c @group
@c @exdent Example:
@c 
@c @end group
@c @end example
@sp 1
See also:  @code{ShortValue, Value, ListIndex, SetValue}
@end deffn









@subsection Scroll Bars
The @code{ScrollBar} class represents a control which allows the user to
select a number between two ranges via a linear, visually intuitive
control.  The range defaults to 0 to 100 and may be controlled via
@code{ScrollBarRange}.

The @code{ScrollBar} class is a subclass of @code{Control} and as such
inherits all the functionality associated with that class.  This section
only documents functionality particular to this class.

Standard control arguments are documented in the section entitled
``Standard Control Method Arguments''.










@deffn {Attach} Attach::ScrollBar
@sp 2
@example
@group
r = gAttach(ctl, val);

object  ctl;   /*  control object  */
object  val;   /*  ctl value       */
object  r;     /*  control object  */
@end group
@end example
This method is used to associate an independent Dynace object with the
state associated with a control object.  @code{val} should be an
instance of the @code{ShortInteger} class.  This object will be
automatically updated to reflect the position of the scroll bar.

This object (@code{val}) will never be disposed by WDS, even after
the control or associated dialog are disposed.  Therefore, this
is one way of gaining access to a control's value after the life
of the control.  It is the programmer's responsibility to dispose of
the object when it is no longer needed.
@example
@group
@exdent Example:

object  ctl, val;

val = gNewWithInt(ShortInteger, 0);
gAttach(ctl, val);
@end group
@end example
@sp 1
See also:  @code{Value, ShortValue, SetValue, SetShortValue,}
        @code{ScrollBarRange}
@end deffn










@deffn {Increment} Increment::ScrollBar
@sp 2
@example
@group
r = gIncrement(ctl, val);

object  ctl;   /*  control object  */
int     val;   /*  increment value */
object  r;     /*  control object  */
@end group
@end example
This method is used to increment the position of the scroll bar
@code{val} points from its current position.  @code{val} may
be negative or positive, which will determine the direction of movement.
@example
@group
@exdent Example:

object  ctl;

gIncrement(ctl, 10);
@end group
@end example
@sp 1
See also:  @code{SetShortValue, SetValue, ScrollBarRange}
@end deffn











@deffn {NewCtl} NewCtl::ScrollBar
@sp 2
@example
@group
ctl = mNewCtl(ScrollBar, id);

unsigned  id;   /*  control id      */
object   ctl;   /*  control object  */
@end group
@end example
@code{mNewCtl} is a macro (defined in @code{dynwin.h}) that expands to
a call to @code{vNew}.

This method is used to create a new control object to be identified as
@code{id} (see section ``Standard Control Method Arguments'').  This
method is mainly used internally.  A programmer would more often
use @code{AddControl::Dialog} to create controls and associate them
to a dialog.
@example
@group
@exdent Example:

object  ctl;

ctl = mNewCtl(ScrollBar, MY_SCROLLBAR);
@end group
@end example
@sp 1
See also:  @code{AddControl::Dialog, GetControl::Dialog}
@end deffn






@deffn {ScrollBarRange} ScrollBarRange::ScrollBar
@sp 2
@example
@group
r = gScrollBarRange(ctl, min, max, pginc, lininc)

object  ctl;    /*  control object  */
int     min;    /*  minimum range   */
int     max;    /*  maximum range   */
int     pginc;  /*  page increment  */
int     lineinc;/*  line increment  */
object  r;      /*  control object  */
@end group
@end example
This method is used to control the range and scrolling characteristics
associated with a given scroll bar control.

@code{min} and @code{max} are used to determine the range of the
control.  These are the values which are returned when the control
is at either end of its extremes.  The default values are 0 and
100 respectively.

@code{pginc} is used to control exactly how much the position of the
scroll bar will change when the user clicks the control in such a way
as to cause a page jump in the position.  This is normally done by
clicking in the area on either side of position indicating element of the
control.  The default value for the variable is 10.  It must be less than
the difference between @code{max} and @code{min}, and is normally larger
than @code{lineinc}.

@code{lineinc} us used to control exactly how much the position of the
scroll bar will change when the user clicks the line increment arrows
located on the extreme ends of the control.  The default value of this
control is 2.
@example
@group
@exdent Example:

object  ctl;

gScrollBarRange(ctl, 0, 100, 10, 2);
@end group
@end example
@c @sp 1
@c See also:  @code{}
@end deffn














@deffn {SetShortValue} SetShortValue::ScrollBar
@sp 2
@example
@group
r = gSetShortValue(ctl, val);

object  ctl;    /*  control object  */
int     val;    /*  ctl value       */
object  r;      /*  control object  */
@end group
@end example
This method is used to set the value associated with a control.  It is
often used to set the initial value prior to performing a dialog.
@code{val} should be the value which represents the desired
position of the scroll bar.  It must be between the minimum and
maximum range associated with the control.

Any previously associated value object (associated via @code{SetValue})
will be disposed when the new value is set.
@example
@group
@exdent Example:

object  ctl;

gSetShortValue(ctl, 50);
@end group
@end example
@sp 1
See also:  @code{SetValue, Increment, CtlShortValue::Dialog,}
        @code{ScrollBarRange}
@end deffn













@deffn {SetValue} SetValue::ScrollBar
@sp 2
@example
@group
r = gSetValue(ctl, val);

object  ctl;    /*  control object  */
object  val;    /*  ctl value       */
object  r;      /*  control object  */
@end group
@end example
This method is used to set the value associated with a control.  It is
often used to set the initial value prior to performing a dialog.
@code{val} should be a Dynace object which is an instance of the
@code{ShortInteger} class and initialized to the value desired for the
control.

@code{val} must represent a number between the minimum and maximum range
associated with the scroll bar.

Any previously associated object will be disposed when the new value is set.
Also, @code{val} will automatically be disposed when the control or associated
dialog is disposed.
@example
@group
@exdent Example:

object  ctl;

gSetValue(ctl, gNewWithInt(ShortInteger, 50));
@end group
@end example
@sp 1
See also:  @code{SetShortValue, Increment, Value, ScrollBarRange}
@end deffn













@deffn {ShortValue} ShortValue::ScrollBar
@sp 2
@example
@group
val = gShortValue(ctl);

object  ctl;   /*  control object  */
short   val;   /*  ctl value       */
@end group
@end example
This method is used to obtain a C short integer which represents the
position of the scroll bar.  The returned value will represent the
position of the scroll bar control.  This number will be between the
minimum and maximum range associated with the control via
@code{ScrollBarRange}.
@example
@group
@exdent Example:

object  ctl;
short   val;

val = gShortValue(ctl);
@end group
@end example
@sp 1
See also:  @code{Value, Attach, SetShortValue, CtlShortValue::Dialog}
@end deffn









@deffn {Value} Value::ScrollBar
@sp 2
@example
@group
val = gValue(ctl);

object  ctl;   /*  control object  */
object  val;   /*  ctl value       */
@end group
@end example
This method is used to obtain a Dynace object which represents the value
associated with the control.  The object returned will be an instance of
the @code{ShortInteger} class.  The returned object will represent
the position of the scroll bar control.  This number will be between
the minimum and maximum range associated with the control via
@code{ScrollBarRange}.

This object will be disposed by WDS when the control object or
associated dialog is disposed.
@example
@group
@exdent Example:

object  ctl, val;

val = gValue(ctl);
@end group
@end example
@sp 1
See also:  @code{ShortValue, Attach, SetValue}
@end deffn









@subsection Custom Controls


